% =============================================================
%  Thesis-style Technical Report Draft
%  Author: Fanbo Meng
%  Institution: Chalmers University of Technology, in collaboration with Viscando AB
%  Date: June 2026
%
%  Usage:
%    Standalone: xelatex technical_report.tex
%    In thesis:  % =============================================================
%  Thesis-style Technical Report Draft
%  Author: Fanbo Meng
%  Institution: Chalmers University of Technology, in collaboration with Viscando AB
%  Date: June 2026
%
%  Usage:
%    Standalone: xelatex technical_report.tex
%    In thesis:  % =============================================================
%  Thesis-style Technical Report Draft
%  Author: Fanbo Meng
%  Institution: Chalmers University of Technology, in collaboration with Viscando AB
%  Date: June 2026
%
%  Usage:
%    Standalone: xelatex technical_report.tex
%    In thesis:  % =============================================================
%  Thesis-style Technical Report Draft
%  Author: Fanbo Meng
%  Institution: Chalmers University of Technology, in collaboration with Viscando AB
%  Date: June 2026
%
%  Usage:
%    Standalone: xelatex technical_report.tex
%    In thesis:  \input{thesis_chapters/technical_report}
% =============================================================

\documentclass[12pt,a4paper]{article}

\usepackage[margin=2.5cm]{geometry}
\usepackage{fontspec}
\defaultfontfeatures{Ligatures=TeX}
\usepackage{amsmath,amssymb}
\usepackage{siunitx}
\usepackage{booktabs}
\usepackage{graphicx}
\usepackage{subcaption}
\usepackage[table]{xcolor}
\usepackage{hyperref}
\usepackage{microtype}
\usepackage{parskip}
\usepackage{enumitem}
\usepackage{array}
\usepackage{float}
\usepackage{tabularx}

\hypersetup{colorlinks=true, linkcolor=blue!70!black, urlcolor=blue!70!black, citecolor=green!50!black}
\graphicspath{{figures/}{../00_pose_pipeline/runs/current_2025_ergonomics/scatter/}}
\emergencystretch=2em

\definecolor{goodcolor}{rgb}{0.00,0.38,0.18}
\definecolor{warncolor}{rgb}{0.55,0.27,0.07}
\definecolor{badcolor}{rgb}{0.65,0.10,0.10}
\newcommand{\good}[1]{\textcolor{goodcolor}{\textbf{#1}}}
\newcommand{\warn}[1]{\textcolor{warncolor}{\textbf{#1}}}
\newcommand{\bad}[1]{\textcolor{badcolor}{\textbf{#1}}}
\newcommand{\degr}{\ensuremath{^\circ}}

\title{\textbf{Posture and Activity Detection in Manufacturing Environments\\Using 3D Vision and AI}}
\author{Fanbo Meng\\
\small Chalmers University of Technology\\
\small In collaboration with Viscando AB}
\date{Technical Thesis Draft, June 2026}

\begin{document}
\maketitle
\tableofcontents
\newpage

% =============================================================
\section*{Abstract}
% =============================================================

Automated ergonomic assessment requires pose-estimation systems that are not
only spatially accurate, but also temporally stable and interpretable at the
joint-angle level.
This thesis studies 3D human pose estimation for manufacturing environments
using a stereo camera setup, external reference measurements from an Xsens
motion-capture suit, and external pose outputs from FastSAM3D-based pipelines.
The central question is how traditional stereo geometry compares with
learning-based pose estimation when the final application is ergonomic risk
assessment rather than generic 3D reconstruction.

The study develops and evaluates a sparse keypoint triangulation pipeline
(SKT), a monocular FastSAM3D baseline, a Merge variant provided as an external
hybrid comparison, and a dense stereo SGBM alternative.
The evaluation is organized around three complementary dimensions:
angle-space accuracy, motion-space consistency, and position-space geometry.
On the current single-subject ergonomics recording, FastSAM3D achieves lower
mean angle error than SKT (15.4\degr\ vs.\ 17.6\degr) and stronger K=6
motion-delta agreement with the Xsens-derived reference (mean Pearson 0.64
vs.\ 0.38).
However, SKT retains the clearest metric-scale 3D interpretation because its
depth is obtained from calibrated stereo geometry, while FastSAM3D requires
additional coordinate and scale handling before position-space comparisons are
meaningful.
The results suggest that learning-based monocular models are competitive for
angle and motion tracking, whereas stereo vision remains valuable for scale,
geometry-aware quality checks, and deployment interpretability.

% =============================================================
\section{Introduction}
% =============================================================

\subsection{Motivation}

Ergonomic risk assessment in manufacturing environments depends on reliable
estimates of human posture over time.
Industrial operators repeatedly reach, lift, bend, and interact with equipment,
and these movements can lead to musculoskeletal risk when joint postures remain
unfavourable or when high-risk actions are repeated frequently.
Manual assessment tools such as RULA are useful but labour-intensive, observer
dependent, and difficult to apply continuously in production settings.
This motivates camera-based systems that can estimate body posture
automatically and provide measurements that are meaningful for ergonomic
analysis.

Human pose estimation is a promising sensing component for this task, but the
industrial setting creates requirements that differ from generic pose
benchmarks.
For ergonomic monitoring, it is not sufficient for the skeleton to look
plausible in an image; joint angles must be stable, temporal motion must be
consistent, and the system should remain interpretable when occlusion, dark
clothing, or viewpoint changes occur.
These requirements make the problem a bridge between computer vision,
multi-view geometry, and application-level ergonomic evaluation.

\subsection{Problem Statement}

This thesis focuses on estimating 3D body keypoints and joint angles from
stereo video in a manufacturing-like environment.
The main method is Sparse Keypoint Triangulation (SKT), where 2D pose detections
from synchronized left and right cameras are triangulated into 3D coordinates.
The study compares this geometry-based approach with external FastSAM3D outputs
and with a dense stereo disparity alternative.
The objective is not to claim one universal best pose estimator, but to
understand which method properties matter for ergonomic assessment.

A central difficulty is that no perfect physical ground-truth system is
available for the full recording.
Xsens measurements are therefore used as an external comparison/reference
system.
This choice is useful because Xsens provides a synchronized full-body motion
signal, but it also requires care because IMU calibration and angle definitions
can introduce systematic offsets.
For that reason, the evaluation includes both absolute angle comparisons and
relative motion comparisons, where frame-to-frame angle changes reduce the
influence of static offsets.

\subsection{Research Questions}

The thesis is guided by three research questions:

\begin{enumerate}[label=\textbf{RQ\arabic*:}, leftmargin=1.8cm]
  \item How accurately can a stereo keypoint triangulation pipeline estimate
        ergonomic joint angles in the current manufacturing dataset?
  \item How does SKT compare with a learning-based FastSAM3D pipeline in
        angle-space, motion-space, and position-space evaluation?
  \item Which practical limitations must be addressed before the system can be
        moved toward real-time deployment on GPU hardware?
\end{enumerate}

\subsection{Contributions}

This thesis makes four main contributions:

\begin{itemize}[leftmargin=1.2cm]
  \item A stereo 3D pose pipeline based on 2D keypoint detection, calibrated
        stereo triangulation, geometric quality filtering, and joint-angle
        computation.
  \item A unified evaluation framework that separates angle-space accuracy,
        motion-space consistency, and position-space geometry instead of relying
        on a single metric.
  \item A revised time-alignment and filtering protocol, including automatic
        offset estimation and a consistent smoothing policy for camera-based
        pose signals.
  \item A comparative analysis showing that FastSAM3D is stronger for angle and
        motion agreement on the current dataset, while stereo geometry remains
        important for metric scale and explainable quality control.
\end{itemize}

\subsection{Thesis Structure}

Section~\ref{sec:background} summarizes background and related work.
Section~\ref{sec:data_setup} describes the dataset, stereo calibration, and
reference data.
Section~\ref{sec:methods} presents the pose-estimation methods.
Section~\ref{sec:evaluation} defines the evaluation framework.
Section~\ref{sec:results} reports the experimental results.
Section~\ref{sec:discussion} discusses applicability and limitations, and
Section~\ref{sec:conclusion} concludes the report.

% =============================================================
\section{Background and Related Work}
\label{sec:background}
% =============================================================

\subsection{3D Human Pose Estimation}

3D human pose estimation recovers body keypoints or body-model parameters from
sensor observations.
Monocular approaches estimate 3D pose from a single image or video stream.
They can use 2D keypoints followed by 3D lifting, or directly regress 3D pose
from visual features.
These methods are attractive for deployment because they require only one
camera, and modern learning-based models can exploit strong priors about human
body shape and motion.
Their main limitation is that absolute scale and depth are not directly
observed from a single view, so position-space evaluation often requires
external scale correction or coordinate alignment.

Stereo and multi-view approaches use geometric constraints between cameras to
recover depth.
In a calibrated stereo setup, corresponding 2D observations in the left and
right cameras define viewing rays whose intersection gives a 3D point estimate.
This gives stereo systems a natural metric-scale interpretation, which is
important when physical distances and workspace geometry matter.
However, sparse triangulation is sensitive to 2D keypoint jitter and to
left--right mismatch, while dense disparity methods can fail on low-texture
clothing or visually ambiguous industrial backgrounds.

\subsection{Ergonomic Assessment from Pose}

RULA is a common ergonomic screening tool that maps upper-limb and body
postures to discrete risk categories.
Because RULA depends primarily on joint-angle ranges, a pose-estimation system
for ergonomics should be evaluated in angle space, not only in 3D position
space.
In this thesis, RULA-like bin agreement is used as an application-level
indicator for selected joints, but the full RULA workflow is not treated as a
completed automated scoring system.
This distinction keeps the evaluation focused on the pose-estimation component.

\subsection{Evaluation Dimensions}

Pose-estimation errors can appear differently depending on the evaluation
space.
Position-space metrics such as MPJPE measure 3D joint distance and are useful
for assessing reconstruction geometry.
Angle-space metrics measure whether the body posture is correct for ergonomic
interpretation.
Motion-space metrics measure whether the temporal change of posture is
consistent, which is especially useful when a reference system may contain a
static calibration offset.
This thesis therefore reports all three dimensions instead of treating any
single metric as sufficient.

\subsection{Reference Systems}

IMU-based motion capture systems such as Xsens provide high-frequency body
motion measurements without requiring optical markers.
They are practical for industrial recordings but rely on sensor calibration and
body-model assumptions.
In this project, Xsens is therefore used as an external comparison/reference
system rather than an absolute physical ground truth.
The thesis also distinguishes between Xsens native angles and Xsens-derived
geometric angles computed from 3D segment positions using the same geometric
formula as the vision methods.
The latter representation is used as the main reference for fairer
cross-method angle comparison.

% =============================================================
\section{Data and Experimental Setup}
\label{sec:data_setup}
% =============================================================

\subsection{Stereo Capture System}

The recording system consists of two synchronized cameras mounted as a stereo
pair.
The videos are recorded at a nominal frame rate of 12.5 fps and a resolution of
2048$\times$1536.
The raw videos are physically inverted due to the camera mounting, so the
current dataset is processed with a 180\degr\ rotation before pose estimation.
Left and right streams are paired using hardware frame identifiers and
metadata timestamps.

\subsection{Camera Calibration}

Stereo calibration determines the camera intrinsics, distortion parameters, and
relative pose between the two cameras.
The current calibration uses an asymmetric circular dot grid and a rational
distortion model.
The optimized calibration gives a baseline of approximately 41.27 cm and
sub-pixel reprojection residuals in the calibration images.
The calibration matrices are stored separately and are treated as fixed inputs
for all SKT and SGBM experiments in this report.

\subsection{2025 Ergonomics Dataset}

The current evaluation uses a single-subject ergonomics recording collected in
a controlled manufacturing-like environment.
The subject performs walking, sitting, squatting, reaching, lifting, and
interaction with furniture.
After left--right pairing by hardware frame ID, the usable synchronized camera
timeline contains 2801 frame pairs over 248.8 s, corresponding to an effective
frame rate of approximately 12.5 fps.
The left metadata contains 3015 rows and the right metadata contains 2882 rows,
so some rows are unpaired or outside the synchronized interval.

\begin{table}[H]
  \centering
  \caption{Summary of the synchronized camera timeline used for evaluation.}
  \label{tab:timeline_summary}
  \begin{tabular}{lc}
    \toprule
    Quantity & Value \\
    \midrule
    Synchronized frame pairs & 2801 \\
    Duration & 248.8 s \\
    Median frame interval & 0.0800 s \\
    Effective frame rate & 12.50 fps \\
    Left metadata rows & 3015 \\
    Right metadata rows & 2882 \\
    Rotation applied & 180\degr \\
    \bottomrule
  \end{tabular}
\end{table}

\subsection{Xsens Reference Data and Angle Representations}

The Xsens recording provides full-body motion data on an independent time axis.
Two angle representations are considered.
The first is the native Xsens joint-angle output, which follows the Xsens body
model and its internal angle definitions.
The second is the Xsens-derived geometric angle representation, where joint
angles are recomputed from Xsens 3D positions using the same three-point
formula as the vision pipelines.
The geometric representation is used as the main reference because it reduces
angle-definition mismatch between systems.

The comparison between Xsens native and Xsens-derived geometric angles is also
used as a reference-system sanity check.
For elbow flexion, the two Xsens-derived representations are very close
(about 1--1.5\degr\ MAE in the current evaluation), whereas some shoulder and
hip angles differ more strongly.
This supports the decision to report motion-space metrics in addition to
absolute angle errors.

\subsection{Cross-system Time Synchronization}

The camera timeline and Xsens timeline are aligned using an automatic offset
search.
The search first evaluates a coarse range from 0 to 30 s, then performs a fine
search around the best candidate.
Three scores are computed: a position-space score, an angle-space score, and a
motion-delta score.
The selected offset is the motion-delta optimum when sufficient samples are
available, because the motion evaluation is most sensitive to temporal
misalignment.

For the current dataset, the selected offset is 17.31 s.
The best position, angle, and motion-delta offsets are close to each other
(17.27 s, 17.28 s, and 17.31 s respectively), which indicates that a constant
offset is a reasonable alignment model for this recording.

% =============================================================
\section{Methods}
\label{sec:methods}
% =============================================================

\subsection{Overview}

This thesis compares four pose-estimation paths under a shared evaluation
protocol.
SKT is the main geometry-based stereo method.
FastSAM3D is an external monocular learning-based baseline.
Merge is an external hybrid output that combines information from FastSAM3D
and Viscando data processing.
SGBM is an alternative dense stereo path used to test whether dense disparity
can improve depth estimation at body joints.

\subsection{Sparse Keypoint Triangulation}

SKT estimates 3D body keypoints by combining 2D pose detection with calibrated
stereo geometry.
Given synchronized left and right video frames, the pipeline first applies a
2D human pose detector to each view.
The resulting COCO-format 2D keypoints are paired across cameras, filtered by
confidence and geometric consistency, and triangulated into 3D coordinates
using the calibrated projection matrices.
Joint angles are then computed from triplets of 3D keypoints.

The motivation for SKT is interpretability.
Each 3D point is linked to two observed 2D detections and a known stereo
calibration model, so the system can expose quality signals such as confidence,
epipolar error, disparity, and reprojection error.
This is valuable in industrial deployment because unreliable frames can be
detected rather than silently converted into ergonomic scores.
The main weakness is that each frame is still strongly dependent on 2D
keypoint quality; small 2D jitter can be amplified by triangulation and then
further amplified by joint-angle computation.

\subsection{FastSAM3D}

FastSAM3D is used as a learning-based monocular comparison method.
Instead of explicitly triangulating left--right keypoint pairs, FastSAM3D
produces 3D pose outputs from a single camera stream using a model-based pose
pipeline.
In this thesis, the FastSAM3D outputs are provided as TRC marker files and are
aligned to the common camera timeline before evaluation.

The motivation for including FastSAM3D is that modern learning-based pose
pipelines can implicitly encode body-shape and motion priors.
This can make the output smoother and more anatomically plausible than a
frame-independent sparse stereo reconstruction, even though the method uses
less direct geometric information.
The limitation is that monocular 3D output may not share the same absolute
scale, coordinate convention, or keypoint semantics as the stereo pipeline.
For this reason, FastSAM3D is evaluated primarily in angle and motion space,
while position-space results are interpreted cautiously.

\subsection{Merge Hybrid Output}

The Merge output is included as an external hybrid comparison.
It is evaluated with the same angle and motion metrics as SKT and FastSAM3D.
The purpose is to test whether combining Viscando processing with FastSAM3D
improves agreement with the reference signal.
In the current results, Merge does not improve the headline angle or
motion-space metrics, so it is treated as a diagnostic comparison rather than
as the main thesis method.

\subsection{Dense Stereo SGBM}

The SGBM path estimates depth from dense stereo disparity rather than sparse
keypoint triangulation.
After stereo rectification, Semi-Global Block Matching computes a disparity
map for each frame.
For each 2D keypoint location, the local disparity is converted to depth using
the stereo baseline and rectified focal length.

The motivation for SGBM is to test whether dense depth can provide a more
robust depth estimate than sparse triangulation when 2D keypoint confidence is
low.
However, dense stereo relies on local texture.
In the current dataset, dark clothing and dark furniture create low-texture
regions near important joints, causing disparity values to collapse or become
unreliable.
SGBM is therefore kept as an alternative path and diagnostic experiment, not as
the main pose-estimation method.

% =============================================================
\section{Evaluation Framework}
\label{sec:evaluation}
% =============================================================

\subsection{Angle-space Evaluation}

Angle-space evaluation measures whether the estimated posture is suitable for
ergonomic interpretation.
Eight full-body angles are evaluated: left and right shoulder, elbow, hip, and
knee flexion.
For each angle, the same three-point geometric formula is applied to the
vision methods and to the Xsens-derived reference:
\begin{equation}
  \theta = \arccos\!\left(
    \frac{(\mathbf{p}_A-\mathbf{p}_B)\cdot(\mathbf{p}_C-\mathbf{p}_B)}
         {\|\mathbf{p}_A-\mathbf{p}_B\|\;\|\mathbf{p}_C-\mathbf{p}_B\|}
  \right).
\end{equation}
The main metrics are MAE, median absolute error, bias, valid-pair ratio, and
RULA-like bin agreement for upper-body joints.

\subsection{Motion-space Evaluation}

Motion-space evaluation compares relative angle changes rather than absolute
angles.
For a given time scale $K$, the K-frame delta is defined as
\begin{equation}
  \Delta_K \theta_t = \theta_t - \theta_{t-K}.
\end{equation}
This reduces the influence of static offsets and focuses on whether two
systems describe the same motion.
The current evaluation reports $K \in \{1,6,12,25\}$, with K=6 used as the
main visualization scale because it corresponds to approximately 0.48 s at
12.5 fps.
Metrics include Pearson correlation, Spearman correlation, delta MAE, RMSE,
high-delta counts, and path ratio.

\subsection{Segment-level Evaluation}

Segment-level evaluation summarizes active movement intervals.
For each detected activity segment, range of motion (ROM), angle MAE, RULA-like
agreement, and dynamic time warping (DTW) distance are computed.
Before DTW, each segment is mean-subtracted and L2-normalized so that the
distance emphasizes motion shape rather than absolute offset or amplitude.
This makes DTW complementary to RMSE and MAE: MAE measures amplitude error,
whereas DTW measures whether the shape of the motion is similar under small
temporal deviations.

\subsection{Position-space Evaluation}

Position-space evaluation measures 3D joint distance and geometric consistency.
For stereo methods, calibrated geometry gives metric-scale coordinates, making
MPJPE and pelvis-relative comparisons meaningful after rigid alignment.
For monocular FastSAM3D output, position-space comparison is more difficult
because the output coordinate system and scale do not necessarily match the
stereo or Xsens coordinate frames.
Therefore, position-space results are used mainly to discuss the practical
advantage of stereo geometry and the coordinate/scale limitations of monocular
methods.

\subsection{Filtering Policy}

The revised evaluation uses a system-specific filtering policy.
Xsens signals are interpolated onto the shared timeline but are not smoothed
again, because they are already heavily processed internally.
Camera-based methods are evaluated after a lightweight smoothing policy
defined in the pipeline configuration.
The important methodological principle is that smoothing is applied before
delta computation; otherwise, the delta metric would measure filter artifacts
rather than the motion signal itself.

% =============================================================
\section{Results}
\label{sec:results}
% =============================================================

\subsection{Angle-space Results}

FastSAM3D achieves the lowest mean angle MAE among the evaluated vision
methods on the current dataset.
Across the eight evaluated angles, FastSAM3D obtains 15.4\degr\ mean MAE,
compared with 17.6\degr\ for SKT and 24.2\degr\ for Merge.
The result indicates that the monocular learning-based output is more stable
for angle estimation in this recording, especially for shoulder, hip, and knee
angles.
The elbow remains difficult for all vision methods.

\begin{table}[H]
  \centering
  \caption{Angle-space summary across eight evaluated joint angles. Lower MAE is better.}
  \label{tab:angle_summary}
  \small
  \begin{tabular}{lcccc}
    \toprule
    Method & Mean MAE $\downarrow$ & Median MAE $\downarrow$ &
    Valid ratio $\uparrow$ & Bin agreement $\uparrow$ \\
    \midrule
    SKT & 17.6 & 17.9 & 0.746 & 0.540 \\
    \good{FastSAM3D} & \good{15.4} & \good{14.4} & \good{0.905} & \good{0.553} \\
    Merge & 24.2 & 22.4 & 0.905 & 0.440 \\
    Xsens native check & 4.7 & 3.7 & 0.905 & 0.886 \\
    \bottomrule
  \end{tabular}
\end{table}

\subsection{Motion-space Results}

FastSAM3D also shows stronger motion agreement than SKT at K=6.
The average K=6 Pearson correlation across the eight angles is 0.64 for
FastSAM3D and 0.38 for SKT.
For the left and right elbow, which are the most important joints in the
motion-space analysis, FastSAM3D reaches Pearson correlations of 0.54 and 0.69,
whereas SKT reaches 0.39 and 0.39.
Merge has near-zero correlation in the current run, indicating that the hybrid
output does not preserve the same short-term motion pattern.

\begin{table}[H]
  \centering
  \caption{K=6 motion-delta agreement. Higher correlation is better.}
  \label{tab:motion_summary}
  \small
  \begin{tabular}{lcccc}
    \toprule
    Method & Mean $r$ $\uparrow$ & Mean $\rho$ $\uparrow$ &
    L-elbow $r$ $\uparrow$ & R-elbow $r$ $\uparrow$ \\
    \midrule
    SKT & 0.379 & 0.334 & 0.392 & 0.386 \\
    \good{FastSAM3D} & \good{0.644} & \good{0.635} & \good{0.542} & \good{0.689} \\
    Merge & 0.006 & 0.027 & 0.016 & 0.081 \\
    Xsens native check & 0.965 & 0.936 & 0.988 & 0.993 \\
    \bottomrule
  \end{tabular}
\end{table}

\begin{figure}[H]
  \centering
  \begin{subfigure}[b]{0.48\linewidth}
    \includegraphics[width=\linewidth]{scatter_skt_vs_xsensfair_leftelbow_k6.png}
    \caption{SKT vs.\ reference, left elbow.}
  \end{subfigure}
  \hfill
  \begin{subfigure}[b]{0.48\linewidth}
    \includegraphics[width=\linewidth]{scatter_fastsam3d_vs_xsensfair_leftelbow_k6.png}
    \caption{FastSAM3D vs.\ reference, left elbow.}
  \end{subfigure}
  \caption{K=6 elbow motion-delta scatter plots. FastSAM3D follows the reference trend more consistently, while SKT shows more high-delta jitter and vertical-axis accumulation.}
  \label{fig:k6_elbow_scatter}
\end{figure}

\subsection{Segment-level Results}

Segment-level ROM and DTW results reinforce the frame-level motion findings.
FastSAM3D has the lowest average ROM absolute error and the lowest DTW distance
among the vision methods, while SKT remains more variable across segments.
The Xsens native check gives the expected lower-bound reference, with small
ROM error and very low DTW distance.

\begin{table}[H]
  \centering
  \caption{Segment-level summary over detected active movement intervals. Lower values are better.}
  \label{tab:segment_summary}
  \begin{tabular}{lccc}
    \toprule
    Method & Segments & Mean ROM error (\degr) $\downarrow$ & Mean DTW distance $\downarrow$ \\
    \midrule
    SKT & 39 & 26.9 & 0.0435 \\
    \good{FastSAM3D} & 45 & \good{14.9} & \good{0.0355} \\
    Merge & 45 & 30.0 & 0.0772 \\
    Xsens native check & 45 & 4.8 & 0.0062 \\
    \bottomrule
  \end{tabular}
\end{table}

\subsection{Position-space Results}

Position-space evaluation gives a different interpretation from the angle and
motion results.
SKT achieves approximately 30.6 cm MPJPE after rigid alignment under the
current distance-evaluation protocol.
FastSAM3D has much larger direct distance to the Xsens-derived reference under
the same coordinate assumptions, because its TRC output uses a different
coordinate and scale convention.
When FastSAM3D is aligned directly to SKT over common joints, the mean
inter-method distance is approximately 60.0 cm.

\begin{table}[H]
  \centering
  \caption{Position-space diagnostic results. These values should be interpreted as coordinate/scale diagnostics rather than pure pose-quality rankings.}
  \label{tab:position_summary}
  \small
  \begin{tabularx}{\linewidth}{lccX}
    \toprule
    Comparison & Mean distance & Median distance & Notes \\
    \midrule
    SKT vs.\ reference & 30.6 cm & 20.8 cm & Metric stereo scale after rigid alignment \\
    FastSAM3D vs.\ reference & 300.3 cm & 294.8 cm & Coordinate/scale convention mismatch dominates \\
    FastSAM3D aligned to SKT & 60.0 cm & 36.3 cm & Direct inter-method skeleton comparison \\
    \bottomrule
  \end{tabularx}
\end{table}

\subsection{Cross-method Synthesis}

The three evaluation dimensions lead to a more nuanced conclusion than a
single metric would provide.
FastSAM3D is currently stronger for angle and motion agreement, suggesting that
learned body and motion priors help suppress jitter in this dataset.
SKT is weaker in short-term motion stability, but it remains valuable because
its 3D coordinates are tied to calibrated stereo geometry and because it
provides explicit quality signals.
Merge does not improve the current headline results.
SGBM shows that dense disparity is not automatically better than sparse
triangulation when the relevant body regions lack reliable texture.

% =============================================================
\section{Discussion}
\label{sec:discussion}
% =============================================================

\subsection{Why FastSAM3D Performs Better in Angle and Motion Space}

FastSAM3D likely benefits from learned priors about human body structure and
motion continuity.
Even though it uses monocular input, the model output is less dependent on
left--right keypoint matching and therefore avoids some stereo-specific jitter
failure modes.
This is particularly visible in the K=6 delta scatter plots, where SKT produces
large short-term deltas even when the reference motion is relatively stable.
Those outliers suggest that the main SKT weakness is not only random noise, but
also occasional keypoint-chain instability around the shoulder--elbow--wrist
geometry.

\subsection{Why Stereo Still Matters}

Stereo geometry remains important because it provides physical scale and
interpretable quality checks.
For real industrial deployment, knowing when the system is uncertain can be as
important as obtaining a smooth pose sequence.
SKT can expose epipolar error, reprojection error, disparity, and left--right
confidence consistency, which are directly tied to the camera geometry.
These signals can support quality gating, uncertainty reporting, or fallback
logic in a production system.

\subsection{Limitations}

The current results are based on a single subject and a single controlled
recording.
Therefore, the quantitative ranking between methods should be treated as
evidence for this dataset rather than a general conclusion about all
manufacturing environments.
The motion-space analysis is strongest for elbow motion because that joint was
investigated most deeply during the project.
The evaluation also depends on alignment, filtering, and segment-detection
parameters, although the current pipeline records these settings explicitly.
Finally, Xsens is an external comparison system rather than an absolute
physical ground truth, so absolute errors should be interpreted with this
limitation in mind.

\subsection{Implications for Real-time Deployment}

The real-time deployment path should first focus on making the existing SKT
pipeline efficient on GPU hardware.
The immediate engineering steps are to export the 2D pose model to TensorRT,
use GPU-accelerated video decoding and preprocessing, and keep triangulation
and angle computation lightweight enough for streaming.
If this baseline runs reliably, transformer-based temporal priors or
FastSAM3D-like model outputs can be explored as optional stabilizers.
FastSAM3D's current runtime on GPU hardware can serve as a practical benchmark
for the expected performance level of learning-based pose methods.

% =============================================================
\section{Conclusion and Future Work}
\label{sec:conclusion}
% =============================================================

\subsection{Conclusion}

This thesis draft presents a multi-dimensional evaluation of 3D pose estimation
for manufacturing-oriented ergonomic assessment.
The main finding is that method quality depends strongly on the evaluation
dimension.
FastSAM3D currently provides better angle-space and motion-space agreement on
the evaluated recording, while SKT provides more interpretable metric 3D
geometry and quality signals.
This suggests that future ergonomic monitoring systems may benefit from a
hybrid view: learning-based models for stable body-pose priors, and stereo
geometry for scale, quality control, and deployment transparency.

\subsection{Future Work}

The next step is validation on additional recordings and subjects.
This is necessary to test whether the current FastSAM3D advantage generalizes
and whether SKT failure modes are consistent across operators, clothing,
lighting, and camera distances.
The second step is real-time GPU deployment of the SKT pipeline, followed by
optional temporal-model integration if the computational budget allows.
The third step is to extend the motion-space evaluation beyond the elbow joint
and connect pose-estimation outputs more directly to full ergonomic risk
assessment.

% =============================================================
\appendix
\clearpage
% =============================================================

\section{Implementation and Reproducibility Notes}

The standalone evaluation workflow is implemented in \texttt{00\_pose\_pipeline/}.
For the current dataset, the main reproduction command is:
\begin{footnotesize}
\begin{verbatim}
/opt/anaconda3/envs/pose/bin/python \
  00_pose_pipeline/src/run_pipeline.py \
  --config \
  00_pose_pipeline/configs/current_2025_ergonomics.yaml \
  --stages \
  validate,offset,angle,motion,segment,scatter
\end{verbatim}
\end{footnotesize}

The configuration records dataset-specific choices such as 180\degr\ video
rotation, hardware-frame-ID synchronization, timestamp parsing, camera
calibration path, smoothing settings, and the automatic offset-search range.
Generated outputs are stored under \texttt{00\_pose\_pipeline/runs/}.

\section{Claim-evidence Self-review}

\begin{table}[H]
  \centering
  \caption{Claim-evidence map for the current thesis draft.}
  \label{tab:claim_evidence}
  \small
  \begin{tabularx}{\linewidth}{p{4.0cm}Xp{1.9cm}}
    \toprule
    Claim & Evidence used in this draft & Status \\
    \midrule
    FastSAM3D is stronger than SKT in angle-space on the current dataset.
    & Table~\ref{tab:angle_summary}: 15.4\degr\ mean MAE vs.\ 17.6\degr.
    & Supported \\
    FastSAM3D is stronger than SKT in K=6 motion-space agreement.
    & Table~\ref{tab:motion_summary}: mean Pearson 0.644 vs.\ 0.379; elbow-specific correlations also higher.
    & Supported \\
    SKT retains a position-space and interpretability advantage.
    & Table~\ref{tab:position_summary}; SKT uses calibrated stereo scale and exposes geometric quality signals.
    & Supported \\
    Results generalize to manufacturing environments broadly.
    & Only one controlled single-subject dataset is evaluated.
    & Needs evidence \\
    Full RULA automation is solved.
    & Only joint-angle and RULA-like bin agreement are evaluated.
    & Not claimed \\
    \bottomrule
  \end{tabularx}
\end{table}

\begin{thebibliography}{9}

\bibitem{kabsch1976solution}
W.~Kabsch,
\textit{A solution for the best rotation to relate two sets of vectors},
Acta Crystallographica, Section A, vol.~32, pp.~922--923, 1976.

\bibitem{hirschmuller2007sgbm}
H.~Hirschm{\"u}ller,
\textit{Stereo processing by semiglobal matching and mutual information},
IEEE Transactions on Pattern Analysis and Machine Intelligence,
vol.~30, no.~2, pp.~328--341, 2007.

\end{thebibliography}

\end{document}

% =============================================================

\documentclass[12pt,a4paper]{article}

\usepackage[margin=2.5cm]{geometry}
\usepackage{fontspec}
\defaultfontfeatures{Ligatures=TeX}
\usepackage{amsmath,amssymb}
\usepackage{siunitx}
\usepackage{booktabs}
\usepackage{graphicx}
\usepackage{subcaption}
\usepackage[table]{xcolor}
\usepackage{hyperref}
\usepackage{microtype}
\usepackage{parskip}
\usepackage{enumitem}
\usepackage{array}
\usepackage{float}
\usepackage{tabularx}

\hypersetup{colorlinks=true, linkcolor=blue!70!black, urlcolor=blue!70!black, citecolor=green!50!black}
\graphicspath{{figures/}{../00_pose_pipeline/runs/current_2025_ergonomics/scatter/}}
\emergencystretch=2em

\definecolor{goodcolor}{rgb}{0.00,0.38,0.18}
\definecolor{warncolor}{rgb}{0.55,0.27,0.07}
\definecolor{badcolor}{rgb}{0.65,0.10,0.10}
\newcommand{\good}[1]{\textcolor{goodcolor}{\textbf{#1}}}
\newcommand{\warn}[1]{\textcolor{warncolor}{\textbf{#1}}}
\newcommand{\bad}[1]{\textcolor{badcolor}{\textbf{#1}}}
\newcommand{\degr}{\ensuremath{^\circ}}

\title{\textbf{Posture and Activity Detection in Manufacturing Environments\\Using 3D Vision and AI}}
\author{Fanbo Meng\\
\small Chalmers University of Technology\\
\small In collaboration with Viscando AB}
\date{Technical Thesis Draft, June 2026}

\begin{document}
\maketitle
\tableofcontents
\newpage

% =============================================================
\section*{Abstract}
% =============================================================

Automated ergonomic assessment requires pose-estimation systems that are not
only spatially accurate, but also temporally stable and interpretable at the
joint-angle level.
This thesis studies 3D human pose estimation for manufacturing environments
using a stereo camera setup, external reference measurements from an Xsens
motion-capture suit, and external pose outputs from FastSAM3D-based pipelines.
The central question is how traditional stereo geometry compares with
learning-based pose estimation when the final application is ergonomic risk
assessment rather than generic 3D reconstruction.

The study develops and evaluates a sparse keypoint triangulation pipeline
(SKT), a monocular FastSAM3D baseline, a Merge variant provided as an external
hybrid comparison, and a dense stereo SGBM alternative.
The evaluation is organized around three complementary dimensions:
angle-space accuracy, motion-space consistency, and position-space geometry.
On the current single-subject ergonomics recording, FastSAM3D achieves lower
mean angle error than SKT (15.4\degr\ vs.\ 17.6\degr) and stronger K=6
motion-delta agreement with the Xsens-derived reference (mean Pearson 0.64
vs.\ 0.38).
However, SKT retains the clearest metric-scale 3D interpretation because its
depth is obtained from calibrated stereo geometry, while FastSAM3D requires
additional coordinate and scale handling before position-space comparisons are
meaningful.
The results suggest that learning-based monocular models are competitive for
angle and motion tracking, whereas stereo vision remains valuable for scale,
geometry-aware quality checks, and deployment interpretability.

% =============================================================
\section{Introduction}
% =============================================================

\subsection{Motivation}

Ergonomic risk assessment in manufacturing environments depends on reliable
estimates of human posture over time.
Industrial operators repeatedly reach, lift, bend, and interact with equipment,
and these movements can lead to musculoskeletal risk when joint postures remain
unfavourable or when high-risk actions are repeated frequently.
Manual assessment tools such as RULA are useful but labour-intensive, observer
dependent, and difficult to apply continuously in production settings.
This motivates camera-based systems that can estimate body posture
automatically and provide measurements that are meaningful for ergonomic
analysis.

Human pose estimation is a promising sensing component for this task, but the
industrial setting creates requirements that differ from generic pose
benchmarks.
For ergonomic monitoring, it is not sufficient for the skeleton to look
plausible in an image; joint angles must be stable, temporal motion must be
consistent, and the system should remain interpretable when occlusion, dark
clothing, or viewpoint changes occur.
These requirements make the problem a bridge between computer vision,
multi-view geometry, and application-level ergonomic evaluation.

\subsection{Problem Statement}

This thesis focuses on estimating 3D body keypoints and joint angles from
stereo video in a manufacturing-like environment.
The main method is Sparse Keypoint Triangulation (SKT), where 2D pose detections
from synchronized left and right cameras are triangulated into 3D coordinates.
The study compares this geometry-based approach with external FastSAM3D outputs
and with a dense stereo disparity alternative.
The objective is not to claim one universal best pose estimator, but to
understand which method properties matter for ergonomic assessment.

A central difficulty is that no perfect physical ground-truth system is
available for the full recording.
Xsens measurements are therefore used as an external comparison/reference
system.
This choice is useful because Xsens provides a synchronized full-body motion
signal, but it also requires care because IMU calibration and angle definitions
can introduce systematic offsets.
For that reason, the evaluation includes both absolute angle comparisons and
relative motion comparisons, where frame-to-frame angle changes reduce the
influence of static offsets.

\subsection{Research Questions}

The thesis is guided by three research questions:

\begin{enumerate}[label=\textbf{RQ\arabic*:}, leftmargin=1.8cm]
  \item How accurately can a stereo keypoint triangulation pipeline estimate
        ergonomic joint angles in the current manufacturing dataset?
  \item How does SKT compare with a learning-based FastSAM3D pipeline in
        angle-space, motion-space, and position-space evaluation?
  \item Which practical limitations must be addressed before the system can be
        moved toward real-time deployment on GPU hardware?
\end{enumerate}

\subsection{Contributions}

This thesis makes four main contributions:

\begin{itemize}[leftmargin=1.2cm]
  \item A stereo 3D pose pipeline based on 2D keypoint detection, calibrated
        stereo triangulation, geometric quality filtering, and joint-angle
        computation.
  \item A unified evaluation framework that separates angle-space accuracy,
        motion-space consistency, and position-space geometry instead of relying
        on a single metric.
  \item A revised time-alignment and filtering protocol, including automatic
        offset estimation and a consistent smoothing policy for camera-based
        pose signals.
  \item A comparative analysis showing that FastSAM3D is stronger for angle and
        motion agreement on the current dataset, while stereo geometry remains
        important for metric scale and explainable quality control.
\end{itemize}

\subsection{Thesis Structure}

Section~\ref{sec:background} summarizes background and related work.
Section~\ref{sec:data_setup} describes the dataset, stereo calibration, and
reference data.
Section~\ref{sec:methods} presents the pose-estimation methods.
Section~\ref{sec:evaluation} defines the evaluation framework.
Section~\ref{sec:results} reports the experimental results.
Section~\ref{sec:discussion} discusses applicability and limitations, and
Section~\ref{sec:conclusion} concludes the report.

% =============================================================
\section{Background and Related Work}
\label{sec:background}
% =============================================================

\subsection{3D Human Pose Estimation}

3D human pose estimation recovers body keypoints or body-model parameters from
sensor observations.
Monocular approaches estimate 3D pose from a single image or video stream.
They can use 2D keypoints followed by 3D lifting, or directly regress 3D pose
from visual features.
These methods are attractive for deployment because they require only one
camera, and modern learning-based models can exploit strong priors about human
body shape and motion.
Their main limitation is that absolute scale and depth are not directly
observed from a single view, so position-space evaluation often requires
external scale correction or coordinate alignment.

Stereo and multi-view approaches use geometric constraints between cameras to
recover depth.
In a calibrated stereo setup, corresponding 2D observations in the left and
right cameras define viewing rays whose intersection gives a 3D point estimate.
This gives stereo systems a natural metric-scale interpretation, which is
important when physical distances and workspace geometry matter.
However, sparse triangulation is sensitive to 2D keypoint jitter and to
left--right mismatch, while dense disparity methods can fail on low-texture
clothing or visually ambiguous industrial backgrounds.

\subsection{Ergonomic Assessment from Pose}

RULA is a common ergonomic screening tool that maps upper-limb and body
postures to discrete risk categories.
Because RULA depends primarily on joint-angle ranges, a pose-estimation system
for ergonomics should be evaluated in angle space, not only in 3D position
space.
In this thesis, RULA-like bin agreement is used as an application-level
indicator for selected joints, but the full RULA workflow is not treated as a
completed automated scoring system.
This distinction keeps the evaluation focused on the pose-estimation component.

\subsection{Evaluation Dimensions}

Pose-estimation errors can appear differently depending on the evaluation
space.
Position-space metrics such as MPJPE measure 3D joint distance and are useful
for assessing reconstruction geometry.
Angle-space metrics measure whether the body posture is correct for ergonomic
interpretation.
Motion-space metrics measure whether the temporal change of posture is
consistent, which is especially useful when a reference system may contain a
static calibration offset.
This thesis therefore reports all three dimensions instead of treating any
single metric as sufficient.

\subsection{Reference Systems}

IMU-based motion capture systems such as Xsens provide high-frequency body
motion measurements without requiring optical markers.
They are practical for industrial recordings but rely on sensor calibration and
body-model assumptions.
In this project, Xsens is therefore used as an external comparison/reference
system rather than an absolute physical ground truth.
The thesis also distinguishes between Xsens native angles and Xsens-derived
geometric angles computed from 3D segment positions using the same geometric
formula as the vision methods.
The latter representation is used as the main reference for fairer
cross-method angle comparison.

% =============================================================
\section{Data and Experimental Setup}
\label{sec:data_setup}
% =============================================================

\subsection{Stereo Capture System}

The recording system consists of two synchronized cameras mounted as a stereo
pair.
The videos are recorded at a nominal frame rate of 12.5 fps and a resolution of
2048$\times$1536.
The raw videos are physically inverted due to the camera mounting, so the
current dataset is processed with a 180\degr\ rotation before pose estimation.
Left and right streams are paired using hardware frame identifiers and
metadata timestamps.

\subsection{Camera Calibration}

Stereo calibration determines the camera intrinsics, distortion parameters, and
relative pose between the two cameras.
The current calibration uses an asymmetric circular dot grid and a rational
distortion model.
The optimized calibration gives a baseline of approximately 41.27 cm and
sub-pixel reprojection residuals in the calibration images.
The calibration matrices are stored separately and are treated as fixed inputs
for all SKT and SGBM experiments in this report.

\subsection{2025 Ergonomics Dataset}

The current evaluation uses a single-subject ergonomics recording collected in
a controlled manufacturing-like environment.
The subject performs walking, sitting, squatting, reaching, lifting, and
interaction with furniture.
After left--right pairing by hardware frame ID, the usable synchronized camera
timeline contains 2801 frame pairs over 248.8 s, corresponding to an effective
frame rate of approximately 12.5 fps.
The left metadata contains 3015 rows and the right metadata contains 2882 rows,
so some rows are unpaired or outside the synchronized interval.

\begin{table}[H]
  \centering
  \caption{Summary of the synchronized camera timeline used for evaluation.}
  \label{tab:timeline_summary}
  \begin{tabular}{lc}
    \toprule
    Quantity & Value \\
    \midrule
    Synchronized frame pairs & 2801 \\
    Duration & 248.8 s \\
    Median frame interval & 0.0800 s \\
    Effective frame rate & 12.50 fps \\
    Left metadata rows & 3015 \\
    Right metadata rows & 2882 \\
    Rotation applied & 180\degr \\
    \bottomrule
  \end{tabular}
\end{table}

\subsection{Xsens Reference Data and Angle Representations}

The Xsens recording provides full-body motion data on an independent time axis.
Two angle representations are considered.
The first is the native Xsens joint-angle output, which follows the Xsens body
model and its internal angle definitions.
The second is the Xsens-derived geometric angle representation, where joint
angles are recomputed from Xsens 3D positions using the same three-point
formula as the vision pipelines.
The geometric representation is used as the main reference because it reduces
angle-definition mismatch between systems.

The comparison between Xsens native and Xsens-derived geometric angles is also
used as a reference-system sanity check.
For elbow flexion, the two Xsens-derived representations are very close
(about 1--1.5\degr\ MAE in the current evaluation), whereas some shoulder and
hip angles differ more strongly.
This supports the decision to report motion-space metrics in addition to
absolute angle errors.

\subsection{Cross-system Time Synchronization}

The camera timeline and Xsens timeline are aligned using an automatic offset
search.
The search first evaluates a coarse range from 0 to 30 s, then performs a fine
search around the best candidate.
Three scores are computed: a position-space score, an angle-space score, and a
motion-delta score.
The selected offset is the motion-delta optimum when sufficient samples are
available, because the motion evaluation is most sensitive to temporal
misalignment.

For the current dataset, the selected offset is 17.31 s.
The best position, angle, and motion-delta offsets are close to each other
(17.27 s, 17.28 s, and 17.31 s respectively), which indicates that a constant
offset is a reasonable alignment model for this recording.

% =============================================================
\section{Methods}
\label{sec:methods}
% =============================================================

\subsection{Overview}

This thesis compares four pose-estimation paths under a shared evaluation
protocol.
SKT is the main geometry-based stereo method.
FastSAM3D is an external monocular learning-based baseline.
Merge is an external hybrid output that combines information from FastSAM3D
and Viscando data processing.
SGBM is an alternative dense stereo path used to test whether dense disparity
can improve depth estimation at body joints.

\subsection{Sparse Keypoint Triangulation}

SKT estimates 3D body keypoints by combining 2D pose detection with calibrated
stereo geometry.
Given synchronized left and right video frames, the pipeline first applies a
2D human pose detector to each view.
The resulting COCO-format 2D keypoints are paired across cameras, filtered by
confidence and geometric consistency, and triangulated into 3D coordinates
using the calibrated projection matrices.
Joint angles are then computed from triplets of 3D keypoints.

The motivation for SKT is interpretability.
Each 3D point is linked to two observed 2D detections and a known stereo
calibration model, so the system can expose quality signals such as confidence,
epipolar error, disparity, and reprojection error.
This is valuable in industrial deployment because unreliable frames can be
detected rather than silently converted into ergonomic scores.
The main weakness is that each frame is still strongly dependent on 2D
keypoint quality; small 2D jitter can be amplified by triangulation and then
further amplified by joint-angle computation.

\subsection{FastSAM3D}

FastSAM3D is used as a learning-based monocular comparison method.
Instead of explicitly triangulating left--right keypoint pairs, FastSAM3D
produces 3D pose outputs from a single camera stream using a model-based pose
pipeline.
In this thesis, the FastSAM3D outputs are provided as TRC marker files and are
aligned to the common camera timeline before evaluation.

The motivation for including FastSAM3D is that modern learning-based pose
pipelines can implicitly encode body-shape and motion priors.
This can make the output smoother and more anatomically plausible than a
frame-independent sparse stereo reconstruction, even though the method uses
less direct geometric information.
The limitation is that monocular 3D output may not share the same absolute
scale, coordinate convention, or keypoint semantics as the stereo pipeline.
For this reason, FastSAM3D is evaluated primarily in angle and motion space,
while position-space results are interpreted cautiously.

\subsection{Merge Hybrid Output}

The Merge output is included as an external hybrid comparison.
It is evaluated with the same angle and motion metrics as SKT and FastSAM3D.
The purpose is to test whether combining Viscando processing with FastSAM3D
improves agreement with the reference signal.
In the current results, Merge does not improve the headline angle or
motion-space metrics, so it is treated as a diagnostic comparison rather than
as the main thesis method.

\subsection{Dense Stereo SGBM}

The SGBM path estimates depth from dense stereo disparity rather than sparse
keypoint triangulation.
After stereo rectification, Semi-Global Block Matching computes a disparity
map for each frame.
For each 2D keypoint location, the local disparity is converted to depth using
the stereo baseline and rectified focal length.

The motivation for SGBM is to test whether dense depth can provide a more
robust depth estimate than sparse triangulation when 2D keypoint confidence is
low.
However, dense stereo relies on local texture.
In the current dataset, dark clothing and dark furniture create low-texture
regions near important joints, causing disparity values to collapse or become
unreliable.
SGBM is therefore kept as an alternative path and diagnostic experiment, not as
the main pose-estimation method.

% =============================================================
\section{Evaluation Framework}
\label{sec:evaluation}
% =============================================================

\subsection{Angle-space Evaluation}

Angle-space evaluation measures whether the estimated posture is suitable for
ergonomic interpretation.
Eight full-body angles are evaluated: left and right shoulder, elbow, hip, and
knee flexion.
For each angle, the same three-point geometric formula is applied to the
vision methods and to the Xsens-derived reference:
\begin{equation}
  \theta = \arccos\!\left(
    \frac{(\mathbf{p}_A-\mathbf{p}_B)\cdot(\mathbf{p}_C-\mathbf{p}_B)}
         {\|\mathbf{p}_A-\mathbf{p}_B\|\;\|\mathbf{p}_C-\mathbf{p}_B\|}
  \right).
\end{equation}
The main metrics are MAE, median absolute error, bias, valid-pair ratio, and
RULA-like bin agreement for upper-body joints.

\subsection{Motion-space Evaluation}

Motion-space evaluation compares relative angle changes rather than absolute
angles.
For a given time scale $K$, the K-frame delta is defined as
\begin{equation}
  \Delta_K \theta_t = \theta_t - \theta_{t-K}.
\end{equation}
This reduces the influence of static offsets and focuses on whether two
systems describe the same motion.
The current evaluation reports $K \in \{1,6,12,25\}$, with K=6 used as the
main visualization scale because it corresponds to approximately 0.48 s at
12.5 fps.
Metrics include Pearson correlation, Spearman correlation, delta MAE, RMSE,
high-delta counts, and path ratio.

\subsection{Segment-level Evaluation}

Segment-level evaluation summarizes active movement intervals.
For each detected activity segment, range of motion (ROM), angle MAE, RULA-like
agreement, and dynamic time warping (DTW) distance are computed.
Before DTW, each segment is mean-subtracted and L2-normalized so that the
distance emphasizes motion shape rather than absolute offset or amplitude.
This makes DTW complementary to RMSE and MAE: MAE measures amplitude error,
whereas DTW measures whether the shape of the motion is similar under small
temporal deviations.

\subsection{Position-space Evaluation}

Position-space evaluation measures 3D joint distance and geometric consistency.
For stereo methods, calibrated geometry gives metric-scale coordinates, making
MPJPE and pelvis-relative comparisons meaningful after rigid alignment.
For monocular FastSAM3D output, position-space comparison is more difficult
because the output coordinate system and scale do not necessarily match the
stereo or Xsens coordinate frames.
Therefore, position-space results are used mainly to discuss the practical
advantage of stereo geometry and the coordinate/scale limitations of monocular
methods.

\subsection{Filtering Policy}

The revised evaluation uses a system-specific filtering policy.
Xsens signals are interpolated onto the shared timeline but are not smoothed
again, because they are already heavily processed internally.
Camera-based methods are evaluated after a lightweight smoothing policy
defined in the pipeline configuration.
The important methodological principle is that smoothing is applied before
delta computation; otherwise, the delta metric would measure filter artifacts
rather than the motion signal itself.

% =============================================================
\section{Results}
\label{sec:results}
% =============================================================

\subsection{Angle-space Results}

FastSAM3D achieves the lowest mean angle MAE among the evaluated vision
methods on the current dataset.
Across the eight evaluated angles, FastSAM3D obtains 15.4\degr\ mean MAE,
compared with 17.6\degr\ for SKT and 24.2\degr\ for Merge.
The result indicates that the monocular learning-based output is more stable
for angle estimation in this recording, especially for shoulder, hip, and knee
angles.
The elbow remains difficult for all vision methods.

\begin{table}[H]
  \centering
  \caption{Angle-space summary across eight evaluated joint angles. Lower MAE is better.}
  \label{tab:angle_summary}
  \small
  \begin{tabular}{lcccc}
    \toprule
    Method & Mean MAE $\downarrow$ & Median MAE $\downarrow$ &
    Valid ratio $\uparrow$ & Bin agreement $\uparrow$ \\
    \midrule
    SKT & 17.6 & 17.9 & 0.746 & 0.540 \\
    \good{FastSAM3D} & \good{15.4} & \good{14.4} & \good{0.905} & \good{0.553} \\
    Merge & 24.2 & 22.4 & 0.905 & 0.440 \\
    Xsens native check & 4.7 & 3.7 & 0.905 & 0.886 \\
    \bottomrule
  \end{tabular}
\end{table}

\subsection{Motion-space Results}

FastSAM3D also shows stronger motion agreement than SKT at K=6.
The average K=6 Pearson correlation across the eight angles is 0.64 for
FastSAM3D and 0.38 for SKT.
For the left and right elbow, which are the most important joints in the
motion-space analysis, FastSAM3D reaches Pearson correlations of 0.54 and 0.69,
whereas SKT reaches 0.39 and 0.39.
Merge has near-zero correlation in the current run, indicating that the hybrid
output does not preserve the same short-term motion pattern.

\begin{table}[H]
  \centering
  \caption{K=6 motion-delta agreement. Higher correlation is better.}
  \label{tab:motion_summary}
  \small
  \begin{tabular}{lcccc}
    \toprule
    Method & Mean $r$ $\uparrow$ & Mean $\rho$ $\uparrow$ &
    L-elbow $r$ $\uparrow$ & R-elbow $r$ $\uparrow$ \\
    \midrule
    SKT & 0.379 & 0.334 & 0.392 & 0.386 \\
    \good{FastSAM3D} & \good{0.644} & \good{0.635} & \good{0.542} & \good{0.689} \\
    Merge & 0.006 & 0.027 & 0.016 & 0.081 \\
    Xsens native check & 0.965 & 0.936 & 0.988 & 0.993 \\
    \bottomrule
  \end{tabular}
\end{table}

\begin{figure}[H]
  \centering
  \begin{subfigure}[b]{0.48\linewidth}
    \includegraphics[width=\linewidth]{scatter_skt_vs_xsensfair_leftelbow_k6.png}
    \caption{SKT vs.\ reference, left elbow.}
  \end{subfigure}
  \hfill
  \begin{subfigure}[b]{0.48\linewidth}
    \includegraphics[width=\linewidth]{scatter_fastsam3d_vs_xsensfair_leftelbow_k6.png}
    \caption{FastSAM3D vs.\ reference, left elbow.}
  \end{subfigure}
  \caption{K=6 elbow motion-delta scatter plots. FastSAM3D follows the reference trend more consistently, while SKT shows more high-delta jitter and vertical-axis accumulation.}
  \label{fig:k6_elbow_scatter}
\end{figure}

\subsection{Segment-level Results}

Segment-level ROM and DTW results reinforce the frame-level motion findings.
FastSAM3D has the lowest average ROM absolute error and the lowest DTW distance
among the vision methods, while SKT remains more variable across segments.
The Xsens native check gives the expected lower-bound reference, with small
ROM error and very low DTW distance.

\begin{table}[H]
  \centering
  \caption{Segment-level summary over detected active movement intervals. Lower values are better.}
  \label{tab:segment_summary}
  \begin{tabular}{lccc}
    \toprule
    Method & Segments & Mean ROM error (\degr) $\downarrow$ & Mean DTW distance $\downarrow$ \\
    \midrule
    SKT & 39 & 26.9 & 0.0435 \\
    \good{FastSAM3D} & 45 & \good{14.9} & \good{0.0355} \\
    Merge & 45 & 30.0 & 0.0772 \\
    Xsens native check & 45 & 4.8 & 0.0062 \\
    \bottomrule
  \end{tabular}
\end{table}

\subsection{Position-space Results}

Position-space evaluation gives a different interpretation from the angle and
motion results.
SKT achieves approximately 30.6 cm MPJPE after rigid alignment under the
current distance-evaluation protocol.
FastSAM3D has much larger direct distance to the Xsens-derived reference under
the same coordinate assumptions, because its TRC output uses a different
coordinate and scale convention.
When FastSAM3D is aligned directly to SKT over common joints, the mean
inter-method distance is approximately 60.0 cm.

\begin{table}[H]
  \centering
  \caption{Position-space diagnostic results. These values should be interpreted as coordinate/scale diagnostics rather than pure pose-quality rankings.}
  \label{tab:position_summary}
  \small
  \begin{tabularx}{\linewidth}{lccX}
    \toprule
    Comparison & Mean distance & Median distance & Notes \\
    \midrule
    SKT vs.\ reference & 30.6 cm & 20.8 cm & Metric stereo scale after rigid alignment \\
    FastSAM3D vs.\ reference & 300.3 cm & 294.8 cm & Coordinate/scale convention mismatch dominates \\
    FastSAM3D aligned to SKT & 60.0 cm & 36.3 cm & Direct inter-method skeleton comparison \\
    \bottomrule
  \end{tabularx}
\end{table}

\subsection{Cross-method Synthesis}

The three evaluation dimensions lead to a more nuanced conclusion than a
single metric would provide.
FastSAM3D is currently stronger for angle and motion agreement, suggesting that
learned body and motion priors help suppress jitter in this dataset.
SKT is weaker in short-term motion stability, but it remains valuable because
its 3D coordinates are tied to calibrated stereo geometry and because it
provides explicit quality signals.
Merge does not improve the current headline results.
SGBM shows that dense disparity is not automatically better than sparse
triangulation when the relevant body regions lack reliable texture.

% =============================================================
\section{Discussion}
\label{sec:discussion}
% =============================================================

\subsection{Why FastSAM3D Performs Better in Angle and Motion Space}

FastSAM3D likely benefits from learned priors about human body structure and
motion continuity.
Even though it uses monocular input, the model output is less dependent on
left--right keypoint matching and therefore avoids some stereo-specific jitter
failure modes.
This is particularly visible in the K=6 delta scatter plots, where SKT produces
large short-term deltas even when the reference motion is relatively stable.
Those outliers suggest that the main SKT weakness is not only random noise, but
also occasional keypoint-chain instability around the shoulder--elbow--wrist
geometry.

\subsection{Why Stereo Still Matters}

Stereo geometry remains important because it provides physical scale and
interpretable quality checks.
For real industrial deployment, knowing when the system is uncertain can be as
important as obtaining a smooth pose sequence.
SKT can expose epipolar error, reprojection error, disparity, and left--right
confidence consistency, which are directly tied to the camera geometry.
These signals can support quality gating, uncertainty reporting, or fallback
logic in a production system.

\subsection{Limitations}

The current results are based on a single subject and a single controlled
recording.
Therefore, the quantitative ranking between methods should be treated as
evidence for this dataset rather than a general conclusion about all
manufacturing environments.
The motion-space analysis is strongest for elbow motion because that joint was
investigated most deeply during the project.
The evaluation also depends on alignment, filtering, and segment-detection
parameters, although the current pipeline records these settings explicitly.
Finally, Xsens is an external comparison system rather than an absolute
physical ground truth, so absolute errors should be interpreted with this
limitation in mind.

\subsection{Implications for Real-time Deployment}

The real-time deployment path should first focus on making the existing SKT
pipeline efficient on GPU hardware.
The immediate engineering steps are to export the 2D pose model to TensorRT,
use GPU-accelerated video decoding and preprocessing, and keep triangulation
and angle computation lightweight enough for streaming.
If this baseline runs reliably, transformer-based temporal priors or
FastSAM3D-like model outputs can be explored as optional stabilizers.
FastSAM3D's current runtime on GPU hardware can serve as a practical benchmark
for the expected performance level of learning-based pose methods.

% =============================================================
\section{Conclusion and Future Work}
\label{sec:conclusion}
% =============================================================

\subsection{Conclusion}

This thesis draft presents a multi-dimensional evaluation of 3D pose estimation
for manufacturing-oriented ergonomic assessment.
The main finding is that method quality depends strongly on the evaluation
dimension.
FastSAM3D currently provides better angle-space and motion-space agreement on
the evaluated recording, while SKT provides more interpretable metric 3D
geometry and quality signals.
This suggests that future ergonomic monitoring systems may benefit from a
hybrid view: learning-based models for stable body-pose priors, and stereo
geometry for scale, quality control, and deployment transparency.

\subsection{Future Work}

The next step is validation on additional recordings and subjects.
This is necessary to test whether the current FastSAM3D advantage generalizes
and whether SKT failure modes are consistent across operators, clothing,
lighting, and camera distances.
The second step is real-time GPU deployment of the SKT pipeline, followed by
optional temporal-model integration if the computational budget allows.
The third step is to extend the motion-space evaluation beyond the elbow joint
and connect pose-estimation outputs more directly to full ergonomic risk
assessment.

% =============================================================
\appendix
\clearpage
% =============================================================

\section{Implementation and Reproducibility Notes}

The standalone evaluation workflow is implemented in \texttt{00\_pose\_pipeline/}.
For the current dataset, the main reproduction command is:
\begin{footnotesize}
\begin{verbatim}
/opt/anaconda3/envs/pose/bin/python \
  00_pose_pipeline/src/run_pipeline.py \
  --config \
  00_pose_pipeline/configs/current_2025_ergonomics.yaml \
  --stages \
  validate,offset,angle,motion,segment,scatter
\end{verbatim}
\end{footnotesize}

The configuration records dataset-specific choices such as 180\degr\ video
rotation, hardware-frame-ID synchronization, timestamp parsing, camera
calibration path, smoothing settings, and the automatic offset-search range.
Generated outputs are stored under \texttt{00\_pose\_pipeline/runs/}.

\section{Claim-evidence Self-review}

\begin{table}[H]
  \centering
  \caption{Claim-evidence map for the current thesis draft.}
  \label{tab:claim_evidence}
  \small
  \begin{tabularx}{\linewidth}{p{4.0cm}Xp{1.9cm}}
    \toprule
    Claim & Evidence used in this draft & Status \\
    \midrule
    FastSAM3D is stronger than SKT in angle-space on the current dataset.
    & Table~\ref{tab:angle_summary}: 15.4\degr\ mean MAE vs.\ 17.6\degr.
    & Supported \\
    FastSAM3D is stronger than SKT in K=6 motion-space agreement.
    & Table~\ref{tab:motion_summary}: mean Pearson 0.644 vs.\ 0.379; elbow-specific correlations also higher.
    & Supported \\
    SKT retains a position-space and interpretability advantage.
    & Table~\ref{tab:position_summary}; SKT uses calibrated stereo scale and exposes geometric quality signals.
    & Supported \\
    Results generalize to manufacturing environments broadly.
    & Only one controlled single-subject dataset is evaluated.
    & Needs evidence \\
    Full RULA automation is solved.
    & Only joint-angle and RULA-like bin agreement are evaluated.
    & Not claimed \\
    \bottomrule
  \end{tabularx}
\end{table}

\begin{thebibliography}{9}

\bibitem{kabsch1976solution}
W.~Kabsch,
\textit{A solution for the best rotation to relate two sets of vectors},
Acta Crystallographica, Section A, vol.~32, pp.~922--923, 1976.

\bibitem{hirschmuller2007sgbm}
H.~Hirschm{\"u}ller,
\textit{Stereo processing by semiglobal matching and mutual information},
IEEE Transactions on Pattern Analysis and Machine Intelligence,
vol.~30, no.~2, pp.~328--341, 2007.

\end{thebibliography}

\end{document}

% =============================================================

\documentclass[12pt,a4paper]{article}

\usepackage[margin=2.5cm]{geometry}
\usepackage{fontspec}
\defaultfontfeatures{Ligatures=TeX}
\usepackage{amsmath,amssymb}
\usepackage{siunitx}
\usepackage{booktabs}
\usepackage{graphicx}
\usepackage{subcaption}
\usepackage[table]{xcolor}
\usepackage{hyperref}
\usepackage{microtype}
\usepackage{parskip}
\usepackage{enumitem}
\usepackage{array}
\usepackage{float}
\usepackage{tabularx}

\hypersetup{colorlinks=true, linkcolor=blue!70!black, urlcolor=blue!70!black, citecolor=green!50!black}
\graphicspath{{figures/}{../00_pose_pipeline/runs/current_2025_ergonomics/scatter/}}
\emergencystretch=2em

\definecolor{goodcolor}{rgb}{0.00,0.38,0.18}
\definecolor{warncolor}{rgb}{0.55,0.27,0.07}
\definecolor{badcolor}{rgb}{0.65,0.10,0.10}
\newcommand{\good}[1]{\textcolor{goodcolor}{\textbf{#1}}}
\newcommand{\warn}[1]{\textcolor{warncolor}{\textbf{#1}}}
\newcommand{\bad}[1]{\textcolor{badcolor}{\textbf{#1}}}
\newcommand{\degr}{\ensuremath{^\circ}}

\title{\textbf{Posture and Activity Detection in Manufacturing Environments\\Using 3D Vision and AI}}
\author{Fanbo Meng\\
\small Chalmers University of Technology\\
\small In collaboration with Viscando AB}
\date{Technical Thesis Draft, June 2026}

\begin{document}
\maketitle
\tableofcontents
\newpage

% =============================================================
\section*{Abstract}
% =============================================================

Automated ergonomic assessment requires pose-estimation systems that are not
only spatially accurate, but also temporally stable and interpretable at the
joint-angle level.
This thesis studies 3D human pose estimation for manufacturing environments
using a stereo camera setup, external reference measurements from an Xsens
motion-capture suit, and external pose outputs from FastSAM3D-based pipelines.
The central question is how traditional stereo geometry compares with
learning-based pose estimation when the final application is ergonomic risk
assessment rather than generic 3D reconstruction.

The study develops and evaluates a sparse keypoint triangulation pipeline
(SKT), a monocular FastSAM3D baseline, a Merge variant provided as an external
hybrid comparison, and a dense stereo SGBM alternative.
The evaluation is organized around three complementary dimensions:
angle-space accuracy, motion-space consistency, and position-space geometry.
On the current single-subject ergonomics recording, FastSAM3D achieves lower
mean angle error than SKT (15.4\degr\ vs.\ 17.6\degr) and stronger K=6
motion-delta agreement with the Xsens-derived reference (mean Pearson 0.64
vs.\ 0.38).
However, SKT retains the clearest metric-scale 3D interpretation because its
depth is obtained from calibrated stereo geometry, while FastSAM3D requires
additional coordinate and scale handling before position-space comparisons are
meaningful.
The results suggest that learning-based monocular models are competitive for
angle and motion tracking, whereas stereo vision remains valuable for scale,
geometry-aware quality checks, and deployment interpretability.

% =============================================================
\section{Introduction}
% =============================================================

\subsection{Motivation}

Ergonomic risk assessment in manufacturing environments depends on reliable
estimates of human posture over time.
Industrial operators repeatedly reach, lift, bend, and interact with equipment,
and these movements can lead to musculoskeletal risk when joint postures remain
unfavourable or when high-risk actions are repeated frequently.
Manual assessment tools such as RULA are useful but labour-intensive, observer
dependent, and difficult to apply continuously in production settings.
This motivates camera-based systems that can estimate body posture
automatically and provide measurements that are meaningful for ergonomic
analysis.

Human pose estimation is a promising sensing component for this task, but the
industrial setting creates requirements that differ from generic pose
benchmarks.
For ergonomic monitoring, it is not sufficient for the skeleton to look
plausible in an image; joint angles must be stable, temporal motion must be
consistent, and the system should remain interpretable when occlusion, dark
clothing, or viewpoint changes occur.
These requirements make the problem a bridge between computer vision,
multi-view geometry, and application-level ergonomic evaluation.

\subsection{Problem Statement}

This thesis focuses on estimating 3D body keypoints and joint angles from
stereo video in a manufacturing-like environment.
The main method is Sparse Keypoint Triangulation (SKT), where 2D pose detections
from synchronized left and right cameras are triangulated into 3D coordinates.
The study compares this geometry-based approach with external FastSAM3D outputs
and with a dense stereo disparity alternative.
The objective is not to claim one universal best pose estimator, but to
understand which method properties matter for ergonomic assessment.

A central difficulty is that no perfect physical ground-truth system is
available for the full recording.
Xsens measurements are therefore used as an external comparison/reference
system.
This choice is useful because Xsens provides a synchronized full-body motion
signal, but it also requires care because IMU calibration and angle definitions
can introduce systematic offsets.
For that reason, the evaluation includes both absolute angle comparisons and
relative motion comparisons, where frame-to-frame angle changes reduce the
influence of static offsets.

\subsection{Research Questions}

The thesis is guided by three research questions:

\begin{enumerate}[label=\textbf{RQ\arabic*:}, leftmargin=1.8cm]
  \item How accurately can a stereo keypoint triangulation pipeline estimate
        ergonomic joint angles in the current manufacturing dataset?
  \item How does SKT compare with a learning-based FastSAM3D pipeline in
        angle-space, motion-space, and position-space evaluation?
  \item Which practical limitations must be addressed before the system can be
        moved toward real-time deployment on GPU hardware?
\end{enumerate}

\subsection{Contributions}

This thesis makes four main contributions:

\begin{itemize}[leftmargin=1.2cm]
  \item A stereo 3D pose pipeline based on 2D keypoint detection, calibrated
        stereo triangulation, geometric quality filtering, and joint-angle
        computation.
  \item A unified evaluation framework that separates angle-space accuracy,
        motion-space consistency, and position-space geometry instead of relying
        on a single metric.
  \item A revised time-alignment and filtering protocol, including automatic
        offset estimation and a consistent smoothing policy for camera-based
        pose signals.
  \item A comparative analysis showing that FastSAM3D is stronger for angle and
        motion agreement on the current dataset, while stereo geometry remains
        important for metric scale and explainable quality control.
\end{itemize}

\subsection{Thesis Structure}

Section~\ref{sec:background} summarizes background and related work.
Section~\ref{sec:data_setup} describes the dataset, stereo calibration, and
reference data.
Section~\ref{sec:methods} presents the pose-estimation methods.
Section~\ref{sec:evaluation} defines the evaluation framework.
Section~\ref{sec:results} reports the experimental results.
Section~\ref{sec:discussion} discusses applicability and limitations, and
Section~\ref{sec:conclusion} concludes the report.

% =============================================================
\section{Background and Related Work}
\label{sec:background}
% =============================================================

\subsection{3D Human Pose Estimation}

3D human pose estimation recovers body keypoints or body-model parameters from
sensor observations.
Monocular approaches estimate 3D pose from a single image or video stream.
They can use 2D keypoints followed by 3D lifting, or directly regress 3D pose
from visual features.
These methods are attractive for deployment because they require only one
camera, and modern learning-based models can exploit strong priors about human
body shape and motion.
Their main limitation is that absolute scale and depth are not directly
observed from a single view, so position-space evaluation often requires
external scale correction or coordinate alignment.

Stereo and multi-view approaches use geometric constraints between cameras to
recover depth.
In a calibrated stereo setup, corresponding 2D observations in the left and
right cameras define viewing rays whose intersection gives a 3D point estimate.
This gives stereo systems a natural metric-scale interpretation, which is
important when physical distances and workspace geometry matter.
However, sparse triangulation is sensitive to 2D keypoint jitter and to
left--right mismatch, while dense disparity methods can fail on low-texture
clothing or visually ambiguous industrial backgrounds.

\subsection{Ergonomic Assessment from Pose}

RULA is a common ergonomic screening tool that maps upper-limb and body
postures to discrete risk categories.
Because RULA depends primarily on joint-angle ranges, a pose-estimation system
for ergonomics should be evaluated in angle space, not only in 3D position
space.
In this thesis, RULA-like bin agreement is used as an application-level
indicator for selected joints, but the full RULA workflow is not treated as a
completed automated scoring system.
This distinction keeps the evaluation focused on the pose-estimation component.

\subsection{Evaluation Dimensions}

Pose-estimation errors can appear differently depending on the evaluation
space.
Position-space metrics such as MPJPE measure 3D joint distance and are useful
for assessing reconstruction geometry.
Angle-space metrics measure whether the body posture is correct for ergonomic
interpretation.
Motion-space metrics measure whether the temporal change of posture is
consistent, which is especially useful when a reference system may contain a
static calibration offset.
This thesis therefore reports all three dimensions instead of treating any
single metric as sufficient.

\subsection{Reference Systems}

IMU-based motion capture systems such as Xsens provide high-frequency body
motion measurements without requiring optical markers.
They are practical for industrial recordings but rely on sensor calibration and
body-model assumptions.
In this project, Xsens is therefore used as an external comparison/reference
system rather than an absolute physical ground truth.
The thesis also distinguishes between Xsens native angles and Xsens-derived
geometric angles computed from 3D segment positions using the same geometric
formula as the vision methods.
The latter representation is used as the main reference for fairer
cross-method angle comparison.

% =============================================================
\section{Data and Experimental Setup}
\label{sec:data_setup}
% =============================================================

\subsection{Stereo Capture System}

The recording system consists of two synchronized cameras mounted as a stereo
pair.
The videos are recorded at a nominal frame rate of 12.5 fps and a resolution of
2048$\times$1536.
The raw videos are physically inverted due to the camera mounting, so the
current dataset is processed with a 180\degr\ rotation before pose estimation.
Left and right streams are paired using hardware frame identifiers and
metadata timestamps.

\subsection{Camera Calibration}

Stereo calibration determines the camera intrinsics, distortion parameters, and
relative pose between the two cameras.
The current calibration uses an asymmetric circular dot grid and a rational
distortion model.
The optimized calibration gives a baseline of approximately 41.27 cm and
sub-pixel reprojection residuals in the calibration images.
The calibration matrices are stored separately and are treated as fixed inputs
for all SKT and SGBM experiments in this report.

\subsection{2025 Ergonomics Dataset}

The current evaluation uses a single-subject ergonomics recording collected in
a controlled manufacturing-like environment.
The subject performs walking, sitting, squatting, reaching, lifting, and
interaction with furniture.
After left--right pairing by hardware frame ID, the usable synchronized camera
timeline contains 2801 frame pairs over 248.8 s, corresponding to an effective
frame rate of approximately 12.5 fps.
The left metadata contains 3015 rows and the right metadata contains 2882 rows,
so some rows are unpaired or outside the synchronized interval.

\begin{table}[H]
  \centering
  \caption{Summary of the synchronized camera timeline used for evaluation.}
  \label{tab:timeline_summary}
  \begin{tabular}{lc}
    \toprule
    Quantity & Value \\
    \midrule
    Synchronized frame pairs & 2801 \\
    Duration & 248.8 s \\
    Median frame interval & 0.0800 s \\
    Effective frame rate & 12.50 fps \\
    Left metadata rows & 3015 \\
    Right metadata rows & 2882 \\
    Rotation applied & 180\degr \\
    \bottomrule
  \end{tabular}
\end{table}

\subsection{Xsens Reference Data and Angle Representations}

The Xsens recording provides full-body motion data on an independent time axis.
Two angle representations are considered.
The first is the native Xsens joint-angle output, which follows the Xsens body
model and its internal angle definitions.
The second is the Xsens-derived geometric angle representation, where joint
angles are recomputed from Xsens 3D positions using the same three-point
formula as the vision pipelines.
The geometric representation is used as the main reference because it reduces
angle-definition mismatch between systems.

The comparison between Xsens native and Xsens-derived geometric angles is also
used as a reference-system sanity check.
For elbow flexion, the two Xsens-derived representations are very close
(about 1--1.5\degr\ MAE in the current evaluation), whereas some shoulder and
hip angles differ more strongly.
This supports the decision to report motion-space metrics in addition to
absolute angle errors.

\subsection{Cross-system Time Synchronization}

The camera timeline and Xsens timeline are aligned using an automatic offset
search.
The search first evaluates a coarse range from 0 to 30 s, then performs a fine
search around the best candidate.
Three scores are computed: a position-space score, an angle-space score, and a
motion-delta score.
The selected offset is the motion-delta optimum when sufficient samples are
available, because the motion evaluation is most sensitive to temporal
misalignment.

For the current dataset, the selected offset is 17.31 s.
The best position, angle, and motion-delta offsets are close to each other
(17.27 s, 17.28 s, and 17.31 s respectively), which indicates that a constant
offset is a reasonable alignment model for this recording.

% =============================================================
\section{Methods}
\label{sec:methods}
% =============================================================

\subsection{Overview}

This thesis compares four pose-estimation paths under a shared evaluation
protocol.
SKT is the main geometry-based stereo method.
FastSAM3D is an external monocular learning-based baseline.
Merge is an external hybrid output that combines information from FastSAM3D
and Viscando data processing.
SGBM is an alternative dense stereo path used to test whether dense disparity
can improve depth estimation at body joints.

\subsection{Sparse Keypoint Triangulation}

SKT estimates 3D body keypoints by combining 2D pose detection with calibrated
stereo geometry.
Given synchronized left and right video frames, the pipeline first applies a
2D human pose detector to each view.
The resulting COCO-format 2D keypoints are paired across cameras, filtered by
confidence and geometric consistency, and triangulated into 3D coordinates
using the calibrated projection matrices.
Joint angles are then computed from triplets of 3D keypoints.

The motivation for SKT is interpretability.
Each 3D point is linked to two observed 2D detections and a known stereo
calibration model, so the system can expose quality signals such as confidence,
epipolar error, disparity, and reprojection error.
This is valuable in industrial deployment because unreliable frames can be
detected rather than silently converted into ergonomic scores.
The main weakness is that each frame is still strongly dependent on 2D
keypoint quality; small 2D jitter can be amplified by triangulation and then
further amplified by joint-angle computation.

\subsection{FastSAM3D}

FastSAM3D is used as a learning-based monocular comparison method.
Instead of explicitly triangulating left--right keypoint pairs, FastSAM3D
produces 3D pose outputs from a single camera stream using a model-based pose
pipeline.
In this thesis, the FastSAM3D outputs are provided as TRC marker files and are
aligned to the common camera timeline before evaluation.

The motivation for including FastSAM3D is that modern learning-based pose
pipelines can implicitly encode body-shape and motion priors.
This can make the output smoother and more anatomically plausible than a
frame-independent sparse stereo reconstruction, even though the method uses
less direct geometric information.
The limitation is that monocular 3D output may not share the same absolute
scale, coordinate convention, or keypoint semantics as the stereo pipeline.
For this reason, FastSAM3D is evaluated primarily in angle and motion space,
while position-space results are interpreted cautiously.

\subsection{Merge Hybrid Output}

The Merge output is included as an external hybrid comparison.
It is evaluated with the same angle and motion metrics as SKT and FastSAM3D.
The purpose is to test whether combining Viscando processing with FastSAM3D
improves agreement with the reference signal.
In the current results, Merge does not improve the headline angle or
motion-space metrics, so it is treated as a diagnostic comparison rather than
as the main thesis method.

\subsection{Dense Stereo SGBM}

The SGBM path estimates depth from dense stereo disparity rather than sparse
keypoint triangulation.
After stereo rectification, Semi-Global Block Matching computes a disparity
map for each frame.
For each 2D keypoint location, the local disparity is converted to depth using
the stereo baseline and rectified focal length.

The motivation for SGBM is to test whether dense depth can provide a more
robust depth estimate than sparse triangulation when 2D keypoint confidence is
low.
However, dense stereo relies on local texture.
In the current dataset, dark clothing and dark furniture create low-texture
regions near important joints, causing disparity values to collapse or become
unreliable.
SGBM is therefore kept as an alternative path and diagnostic experiment, not as
the main pose-estimation method.

% =============================================================
\section{Evaluation Framework}
\label{sec:evaluation}
% =============================================================

\subsection{Angle-space Evaluation}

Angle-space evaluation measures whether the estimated posture is suitable for
ergonomic interpretation.
Eight full-body angles are evaluated: left and right shoulder, elbow, hip, and
knee flexion.
For each angle, the same three-point geometric formula is applied to the
vision methods and to the Xsens-derived reference:
\begin{equation}
  \theta = \arccos\!\left(
    \frac{(\mathbf{p}_A-\mathbf{p}_B)\cdot(\mathbf{p}_C-\mathbf{p}_B)}
         {\|\mathbf{p}_A-\mathbf{p}_B\|\;\|\mathbf{p}_C-\mathbf{p}_B\|}
  \right).
\end{equation}
The main metrics are MAE, median absolute error, bias, valid-pair ratio, and
RULA-like bin agreement for upper-body joints.

\subsection{Motion-space Evaluation}

Motion-space evaluation compares relative angle changes rather than absolute
angles.
For a given time scale $K$, the K-frame delta is defined as
\begin{equation}
  \Delta_K \theta_t = \theta_t - \theta_{t-K}.
\end{equation}
This reduces the influence of static offsets and focuses on whether two
systems describe the same motion.
The current evaluation reports $K \in \{1,6,12,25\}$, with K=6 used as the
main visualization scale because it corresponds to approximately 0.48 s at
12.5 fps.
Metrics include Pearson correlation, Spearman correlation, delta MAE, RMSE,
high-delta counts, and path ratio.

\subsection{Segment-level Evaluation}

Segment-level evaluation summarizes active movement intervals.
For each detected activity segment, range of motion (ROM), angle MAE, RULA-like
agreement, and dynamic time warping (DTW) distance are computed.
Before DTW, each segment is mean-subtracted and L2-normalized so that the
distance emphasizes motion shape rather than absolute offset or amplitude.
This makes DTW complementary to RMSE and MAE: MAE measures amplitude error,
whereas DTW measures whether the shape of the motion is similar under small
temporal deviations.

\subsection{Position-space Evaluation}

Position-space evaluation measures 3D joint distance and geometric consistency.
For stereo methods, calibrated geometry gives metric-scale coordinates, making
MPJPE and pelvis-relative comparisons meaningful after rigid alignment.
For monocular FastSAM3D output, position-space comparison is more difficult
because the output coordinate system and scale do not necessarily match the
stereo or Xsens coordinate frames.
Therefore, position-space results are used mainly to discuss the practical
advantage of stereo geometry and the coordinate/scale limitations of monocular
methods.

\subsection{Filtering Policy}

The revised evaluation uses a system-specific filtering policy.
Xsens signals are interpolated onto the shared timeline but are not smoothed
again, because they are already heavily processed internally.
Camera-based methods are evaluated after a lightweight smoothing policy
defined in the pipeline configuration.
The important methodological principle is that smoothing is applied before
delta computation; otherwise, the delta metric would measure filter artifacts
rather than the motion signal itself.

% =============================================================
\section{Results}
\label{sec:results}
% =============================================================

\subsection{Angle-space Results}

FastSAM3D achieves the lowest mean angle MAE among the evaluated vision
methods on the current dataset.
Across the eight evaluated angles, FastSAM3D obtains 15.4\degr\ mean MAE,
compared with 17.6\degr\ for SKT and 24.2\degr\ for Merge.
The result indicates that the monocular learning-based output is more stable
for angle estimation in this recording, especially for shoulder, hip, and knee
angles.
The elbow remains difficult for all vision methods.

\begin{table}[H]
  \centering
  \caption{Angle-space summary across eight evaluated joint angles. Lower MAE is better.}
  \label{tab:angle_summary}
  \small
  \begin{tabular}{lcccc}
    \toprule
    Method & Mean MAE $\downarrow$ & Median MAE $\downarrow$ &
    Valid ratio $\uparrow$ & Bin agreement $\uparrow$ \\
    \midrule
    SKT & 17.6 & 17.9 & 0.746 & 0.540 \\
    \good{FastSAM3D} & \good{15.4} & \good{14.4} & \good{0.905} & \good{0.553} \\
    Merge & 24.2 & 22.4 & 0.905 & 0.440 \\
    Xsens native check & 4.7 & 3.7 & 0.905 & 0.886 \\
    \bottomrule
  \end{tabular}
\end{table}

\subsection{Motion-space Results}

FastSAM3D also shows stronger motion agreement than SKT at K=6.
The average K=6 Pearson correlation across the eight angles is 0.64 for
FastSAM3D and 0.38 for SKT.
For the left and right elbow, which are the most important joints in the
motion-space analysis, FastSAM3D reaches Pearson correlations of 0.54 and 0.69,
whereas SKT reaches 0.39 and 0.39.
Merge has near-zero correlation in the current run, indicating that the hybrid
output does not preserve the same short-term motion pattern.

\begin{table}[H]
  \centering
  \caption{K=6 motion-delta agreement. Higher correlation is better.}
  \label{tab:motion_summary}
  \small
  \begin{tabular}{lcccc}
    \toprule
    Method & Mean $r$ $\uparrow$ & Mean $\rho$ $\uparrow$ &
    L-elbow $r$ $\uparrow$ & R-elbow $r$ $\uparrow$ \\
    \midrule
    SKT & 0.379 & 0.334 & 0.392 & 0.386 \\
    \good{FastSAM3D} & \good{0.644} & \good{0.635} & \good{0.542} & \good{0.689} \\
    Merge & 0.006 & 0.027 & 0.016 & 0.081 \\
    Xsens native check & 0.965 & 0.936 & 0.988 & 0.993 \\
    \bottomrule
  \end{tabular}
\end{table}

\begin{figure}[H]
  \centering
  \begin{subfigure}[b]{0.48\linewidth}
    \includegraphics[width=\linewidth]{scatter_skt_vs_xsensfair_leftelbow_k6.png}
    \caption{SKT vs.\ reference, left elbow.}
  \end{subfigure}
  \hfill
  \begin{subfigure}[b]{0.48\linewidth}
    \includegraphics[width=\linewidth]{scatter_fastsam3d_vs_xsensfair_leftelbow_k6.png}
    \caption{FastSAM3D vs.\ reference, left elbow.}
  \end{subfigure}
  \caption{K=6 elbow motion-delta scatter plots. FastSAM3D follows the reference trend more consistently, while SKT shows more high-delta jitter and vertical-axis accumulation.}
  \label{fig:k6_elbow_scatter}
\end{figure}

\subsection{Segment-level Results}

Segment-level ROM and DTW results reinforce the frame-level motion findings.
FastSAM3D has the lowest average ROM absolute error and the lowest DTW distance
among the vision methods, while SKT remains more variable across segments.
The Xsens native check gives the expected lower-bound reference, with small
ROM error and very low DTW distance.

\begin{table}[H]
  \centering
  \caption{Segment-level summary over detected active movement intervals. Lower values are better.}
  \label{tab:segment_summary}
  \begin{tabular}{lccc}
    \toprule
    Method & Segments & Mean ROM error (\degr) $\downarrow$ & Mean DTW distance $\downarrow$ \\
    \midrule
    SKT & 39 & 26.9 & 0.0435 \\
    \good{FastSAM3D} & 45 & \good{14.9} & \good{0.0355} \\
    Merge & 45 & 30.0 & 0.0772 \\
    Xsens native check & 45 & 4.8 & 0.0062 \\
    \bottomrule
  \end{tabular}
\end{table}

\subsection{Position-space Results}

Position-space evaluation gives a different interpretation from the angle and
motion results.
SKT achieves approximately 30.6 cm MPJPE after rigid alignment under the
current distance-evaluation protocol.
FastSAM3D has much larger direct distance to the Xsens-derived reference under
the same coordinate assumptions, because its TRC output uses a different
coordinate and scale convention.
When FastSAM3D is aligned directly to SKT over common joints, the mean
inter-method distance is approximately 60.0 cm.

\begin{table}[H]
  \centering
  \caption{Position-space diagnostic results. These values should be interpreted as coordinate/scale diagnostics rather than pure pose-quality rankings.}
  \label{tab:position_summary}
  \small
  \begin{tabularx}{\linewidth}{lccX}
    \toprule
    Comparison & Mean distance & Median distance & Notes \\
    \midrule
    SKT vs.\ reference & 30.6 cm & 20.8 cm & Metric stereo scale after rigid alignment \\
    FastSAM3D vs.\ reference & 300.3 cm & 294.8 cm & Coordinate/scale convention mismatch dominates \\
    FastSAM3D aligned to SKT & 60.0 cm & 36.3 cm & Direct inter-method skeleton comparison \\
    \bottomrule
  \end{tabularx}
\end{table}

\subsection{Cross-method Synthesis}

The three evaluation dimensions lead to a more nuanced conclusion than a
single metric would provide.
FastSAM3D is currently stronger for angle and motion agreement, suggesting that
learned body and motion priors help suppress jitter in this dataset.
SKT is weaker in short-term motion stability, but it remains valuable because
its 3D coordinates are tied to calibrated stereo geometry and because it
provides explicit quality signals.
Merge does not improve the current headline results.
SGBM shows that dense disparity is not automatically better than sparse
triangulation when the relevant body regions lack reliable texture.

% =============================================================
\section{Discussion}
\label{sec:discussion}
% =============================================================

\subsection{Why FastSAM3D Performs Better in Angle and Motion Space}

FastSAM3D likely benefits from learned priors about human body structure and
motion continuity.
Even though it uses monocular input, the model output is less dependent on
left--right keypoint matching and therefore avoids some stereo-specific jitter
failure modes.
This is particularly visible in the K=6 delta scatter plots, where SKT produces
large short-term deltas even when the reference motion is relatively stable.
Those outliers suggest that the main SKT weakness is not only random noise, but
also occasional keypoint-chain instability around the shoulder--elbow--wrist
geometry.

\subsection{Why Stereo Still Matters}

Stereo geometry remains important because it provides physical scale and
interpretable quality checks.
For real industrial deployment, knowing when the system is uncertain can be as
important as obtaining a smooth pose sequence.
SKT can expose epipolar error, reprojection error, disparity, and left--right
confidence consistency, which are directly tied to the camera geometry.
These signals can support quality gating, uncertainty reporting, or fallback
logic in a production system.

\subsection{Limitations}

The current results are based on a single subject and a single controlled
recording.
Therefore, the quantitative ranking between methods should be treated as
evidence for this dataset rather than a general conclusion about all
manufacturing environments.
The motion-space analysis is strongest for elbow motion because that joint was
investigated most deeply during the project.
The evaluation also depends on alignment, filtering, and segment-detection
parameters, although the current pipeline records these settings explicitly.
Finally, Xsens is an external comparison system rather than an absolute
physical ground truth, so absolute errors should be interpreted with this
limitation in mind.

\subsection{Implications for Real-time Deployment}

The real-time deployment path should first focus on making the existing SKT
pipeline efficient on GPU hardware.
The immediate engineering steps are to export the 2D pose model to TensorRT,
use GPU-accelerated video decoding and preprocessing, and keep triangulation
and angle computation lightweight enough for streaming.
If this baseline runs reliably, transformer-based temporal priors or
FastSAM3D-like model outputs can be explored as optional stabilizers.
FastSAM3D's current runtime on GPU hardware can serve as a practical benchmark
for the expected performance level of learning-based pose methods.

% =============================================================
\section{Conclusion and Future Work}
\label{sec:conclusion}
% =============================================================

\subsection{Conclusion}

This thesis draft presents a multi-dimensional evaluation of 3D pose estimation
for manufacturing-oriented ergonomic assessment.
The main finding is that method quality depends strongly on the evaluation
dimension.
FastSAM3D currently provides better angle-space and motion-space agreement on
the evaluated recording, while SKT provides more interpretable metric 3D
geometry and quality signals.
This suggests that future ergonomic monitoring systems may benefit from a
hybrid view: learning-based models for stable body-pose priors, and stereo
geometry for scale, quality control, and deployment transparency.

\subsection{Future Work}

The next step is validation on additional recordings and subjects.
This is necessary to test whether the current FastSAM3D advantage generalizes
and whether SKT failure modes are consistent across operators, clothing,
lighting, and camera distances.
The second step is real-time GPU deployment of the SKT pipeline, followed by
optional temporal-model integration if the computational budget allows.
The third step is to extend the motion-space evaluation beyond the elbow joint
and connect pose-estimation outputs more directly to full ergonomic risk
assessment.

% =============================================================
\appendix
\clearpage
% =============================================================

\section{Implementation and Reproducibility Notes}

The standalone evaluation workflow is implemented in \texttt{00\_pose\_pipeline/}.
For the current dataset, the main reproduction command is:
\begin{footnotesize}
\begin{verbatim}
/opt/anaconda3/envs/pose/bin/python \
  00_pose_pipeline/src/run_pipeline.py \
  --config \
  00_pose_pipeline/configs/current_2025_ergonomics.yaml \
  --stages \
  validate,offset,angle,motion,segment,scatter
\end{verbatim}
\end{footnotesize}

The configuration records dataset-specific choices such as 180\degr\ video
rotation, hardware-frame-ID synchronization, timestamp parsing, camera
calibration path, smoothing settings, and the automatic offset-search range.
Generated outputs are stored under \texttt{00\_pose\_pipeline/runs/}.

\section{Claim-evidence Self-review}

\begin{table}[H]
  \centering
  \caption{Claim-evidence map for the current thesis draft.}
  \label{tab:claim_evidence}
  \small
  \begin{tabularx}{\linewidth}{p{4.0cm}Xp{1.9cm}}
    \toprule
    Claim & Evidence used in this draft & Status \\
    \midrule
    FastSAM3D is stronger than SKT in angle-space on the current dataset.
    & Table~\ref{tab:angle_summary}: 15.4\degr\ mean MAE vs.\ 17.6\degr.
    & Supported \\
    FastSAM3D is stronger than SKT in K=6 motion-space agreement.
    & Table~\ref{tab:motion_summary}: mean Pearson 0.644 vs.\ 0.379; elbow-specific correlations also higher.
    & Supported \\
    SKT retains a position-space and interpretability advantage.
    & Table~\ref{tab:position_summary}; SKT uses calibrated stereo scale and exposes geometric quality signals.
    & Supported \\
    Results generalize to manufacturing environments broadly.
    & Only one controlled single-subject dataset is evaluated.
    & Needs evidence \\
    Full RULA automation is solved.
    & Only joint-angle and RULA-like bin agreement are evaluated.
    & Not claimed \\
    \bottomrule
  \end{tabularx}
\end{table}

\begin{thebibliography}{9}

\bibitem{kabsch1976solution}
W.~Kabsch,
\textit{A solution for the best rotation to relate two sets of vectors},
Acta Crystallographica, Section A, vol.~32, pp.~922--923, 1976.

\bibitem{hirschmuller2007sgbm}
H.~Hirschm{\"u}ller,
\textit{Stereo processing by semiglobal matching and mutual information},
IEEE Transactions on Pattern Analysis and Machine Intelligence,
vol.~30, no.~2, pp.~328--341, 2007.

\end{thebibliography}

\end{document}

% =============================================================

\documentclass[12pt,a4paper]{article}

\usepackage[margin=2.5cm]{geometry}
\usepackage{fontspec}
\defaultfontfeatures{Ligatures=TeX}
\usepackage{amsmath,amssymb}
\usepackage{siunitx}
\usepackage{booktabs}
\usepackage{graphicx}
\usepackage{subcaption}
\usepackage[table]{xcolor}
\usepackage{hyperref}
\usepackage{microtype}
\usepackage{parskip}
\usepackage{enumitem}
\usepackage{array}
\usepackage{float}
\usepackage{tabularx}

\hypersetup{colorlinks=true, linkcolor=blue!70!black, urlcolor=blue!70!black, citecolor=green!50!black}
\graphicspath{{figures/}{../00_pose_pipeline/runs/current_2025_ergonomics/scatter/}}
\emergencystretch=2em

\definecolor{goodcolor}{rgb}{0.00,0.38,0.18}
\definecolor{warncolor}{rgb}{0.55,0.27,0.07}
\definecolor{badcolor}{rgb}{0.65,0.10,0.10}
\newcommand{\good}[1]{\textcolor{goodcolor}{\textbf{#1}}}
\newcommand{\warn}[1]{\textcolor{warncolor}{\textbf{#1}}}
\newcommand{\bad}[1]{\textcolor{badcolor}{\textbf{#1}}}
\newcommand{\degr}{\ensuremath{^\circ}}

\title{\textbf{Posture and Activity Detection in Manufacturing Environments\\Using 3D Vision and AI}}
\author{Fanbo Meng\\
\small Chalmers University of Technology\\
\small In collaboration with Viscando AB}
\date{Technical Thesis Draft, June 2026}

\begin{document}
\maketitle
\tableofcontents
\newpage

% =============================================================
\section*{Abstract}
% =============================================================

Automated ergonomic assessment requires pose-estimation systems that are not
only spatially accurate, but also temporally stable and interpretable at the
joint-angle level.
This thesis studies 3D human pose estimation for manufacturing environments
using a stereo camera setup, external reference measurements from an Xsens
motion-capture suit, and external pose outputs from FastSAM3D-based pipelines.
The central question is how traditional stereo geometry compares with
learning-based pose estimation when the final application is ergonomic risk
assessment rather than generic 3D reconstruction.

The study develops and evaluates a sparse keypoint triangulation pipeline
(SKT), a monocular FastSAM3D baseline, a Merge variant provided as an external
hybrid comparison, and a dense stereo SGBM alternative.
The evaluation is organized around three complementary dimensions:
angle-space accuracy, motion-space consistency, and position-space geometry.
On the current single-subject ergonomics recording, FastSAM3D achieves lower
mean angle error than SKT (15.4\degr\ vs.\ 17.6\degr) and stronger K=6
motion-delta agreement with the Xsens-derived reference (mean Pearson 0.64
vs.\ 0.38).
However, SKT retains the clearest metric-scale 3D interpretation because its
depth is obtained from calibrated stereo geometry, while FastSAM3D requires
additional coordinate and scale handling before position-space comparisons are
meaningful.
The results suggest that learning-based monocular models are competitive for
angle and motion tracking, whereas stereo vision remains valuable for scale,
geometry-aware quality checks, and deployment interpretability.

% =============================================================
\section{Introduction}
% =============================================================

\subsection{Motivation}

Ergonomic risk assessment in manufacturing environments depends on reliable
estimates of human posture over time.
Industrial operators repeatedly reach, lift, bend, and interact with equipment,
and these movements can lead to musculoskeletal risk when joint postures remain
unfavourable or when high-risk actions are repeated frequently.
Manual assessment tools such as RULA are useful but labour-intensive, observer
dependent, and difficult to apply continuously in production settings.
This motivates camera-based systems that can estimate body posture
automatically and provide measurements that are meaningful for ergonomic
analysis.

Human pose estimation is a promising sensing component for this task, but the
industrial setting creates requirements that differ from generic pose
benchmarks.
For ergonomic monitoring, it is not sufficient for the skeleton to look
plausible in an image; joint angles must be stable, temporal motion must be
consistent, and the system should remain interpretable when occlusion, dark
clothing, or viewpoint changes occur.
These requirements make the problem a bridge between computer vision,
multi-view geometry, and application-level ergonomic evaluation.

\subsection{Problem Statement}

This thesis focuses on estimating 3D body keypoints and joint angles from
stereo video in a manufacturing-like environment.
The main method is Sparse Keypoint Triangulation (SKT), where 2D pose detections
from synchronized left and right cameras are triangulated into 3D coordinates.
The study compares this geometry-based approach with external FastSAM3D outputs
and with a dense stereo disparity alternative.
The objective is not to claim one universal best pose estimator, but to
understand which method properties matter for ergonomic assessment.

A central difficulty is that no perfect physical ground-truth system is
available for the full recording.
Xsens measurements are therefore used as an external comparison/reference
system.
This choice is useful because Xsens provides a synchronized full-body motion
signal, but it also requires care because IMU calibration and angle definitions
can introduce systematic offsets.
For that reason, the evaluation includes both absolute angle comparisons and
relative motion comparisons, where frame-to-frame angle changes reduce the
influence of static offsets.

\subsection{Research Questions}

The thesis is guided by three research questions:

\begin{enumerate}[label=\textbf{RQ\arabic*:}, leftmargin=1.8cm]
  \item How accurately can a stereo keypoint triangulation pipeline estimate
        ergonomic joint angles in the current manufacturing dataset?
  \item How does SKT compare with a learning-based FastSAM3D pipeline in
        angle-space, motion-space, and position-space evaluation?
  \item Which practical limitations must be addressed before the system can be
        moved toward real-time deployment on GPU hardware?
\end{enumerate}

\subsection{Contributions}

This thesis makes four main contributions:

\begin{itemize}[leftmargin=1.2cm]
  \item A stereo 3D pose pipeline based on 2D keypoint detection, calibrated
        stereo triangulation, geometric quality filtering, and joint-angle
        computation.
  \item A unified evaluation framework that separates angle-space accuracy,
        motion-space consistency, and position-space geometry instead of relying
        on a single metric.
  \item A revised time-alignment and filtering protocol, including automatic
        offset estimation and a consistent smoothing policy for camera-based
        pose signals.
  \item A comparative analysis showing that FastSAM3D is stronger for angle and
        motion agreement on the current dataset, while stereo geometry remains
        important for metric scale and explainable quality control.
\end{itemize}

\subsection{Thesis Structure}

Section~\ref{sec:background} summarizes background and related work.
Section~\ref{sec:data_setup} describes the dataset, stereo calibration, and
reference data.
Section~\ref{sec:methods} presents the pose-estimation methods.
Section~\ref{sec:evaluation} defines the evaluation framework.
Section~\ref{sec:results} reports the experimental results.
Section~\ref{sec:discussion} discusses applicability and limitations, and
Section~\ref{sec:conclusion} concludes the report.

% =============================================================
\section{Background and Related Work}
\label{sec:background}
% =============================================================

\subsection{3D Human Pose Estimation}

3D human pose estimation recovers body keypoints or body-model parameters from
sensor observations.
Monocular approaches estimate 3D pose from a single image or video stream.
They can use 2D keypoints followed by 3D lifting, or directly regress 3D pose
from visual features.
These methods are attractive for deployment because they require only one
camera, and modern learning-based models can exploit strong priors about human
body shape and motion.
Their main limitation is that absolute scale and depth are not directly
observed from a single view, so position-space evaluation often requires
external scale correction or coordinate alignment.

Stereo and multi-view approaches use geometric constraints between cameras to
recover depth.
In a calibrated stereo setup, corresponding 2D observations in the left and
right cameras define viewing rays whose intersection gives a 3D point estimate.
This gives stereo systems a natural metric-scale interpretation, which is
important when physical distances and workspace geometry matter.
However, sparse triangulation is sensitive to 2D keypoint jitter and to
left--right mismatch, while dense disparity methods can fail on low-texture
clothing or visually ambiguous industrial backgrounds.

\subsection{Ergonomic Assessment from Pose}

RULA is a common ergonomic screening tool that maps upper-limb and body
postures to discrete risk categories.
Because RULA depends primarily on joint-angle ranges, a pose-estimation system
for ergonomics should be evaluated in angle space, not only in 3D position
space.
In this thesis, RULA-like bin agreement is used as an application-level
indicator for selected joints, but the full RULA workflow is not treated as a
completed automated scoring system.
This distinction keeps the evaluation focused on the pose-estimation component.

\subsection{Evaluation Dimensions}

Pose-estimation errors can appear differently depending on the evaluation
space.
Position-space metrics such as MPJPE measure 3D joint distance and are useful
for assessing reconstruction geometry.
Angle-space metrics measure whether the body posture is correct for ergonomic
interpretation.
Motion-space metrics measure whether the temporal change of posture is
consistent, which is especially useful when a reference system may contain a
static calibration offset.
This thesis therefore reports all three dimensions instead of treating any
single metric as sufficient.

\subsection{Reference Systems}

IMU-based motion capture systems such as Xsens provide high-frequency body
motion measurements without requiring optical markers.
They are practical for industrial recordings but rely on sensor calibration and
body-model assumptions.
In this project, Xsens is therefore used as an external comparison/reference
system rather than an absolute physical ground truth.
The thesis also distinguishes between Xsens native angles and Xsens-derived
geometric angles computed from 3D segment positions using the same geometric
formula as the vision methods.
The latter representation is used as the main reference for fairer
cross-method angle comparison.

% =============================================================
\section{Data and Experimental Setup}
\label{sec:data_setup}
% =============================================================

\subsection{Stereo Capture System}

The recording system consists of two synchronized cameras mounted as a stereo
pair.
The videos are recorded at a nominal frame rate of 12.5 fps and a resolution of
2048$\times$1536.
The raw videos are physically inverted due to the camera mounting, so the
current dataset is processed with a 180\degr\ rotation before pose estimation.
Left and right streams are paired using hardware frame identifiers and
metadata timestamps.

\subsection{Camera Calibration}

Stereo calibration determines the camera intrinsics, distortion parameters, and
relative pose between the two cameras.
The current calibration uses an asymmetric circular dot grid and a rational
distortion model.
The optimized calibration gives a baseline of approximately 41.27 cm and
sub-pixel reprojection residuals in the calibration images.
The calibration matrices are stored separately and are treated as fixed inputs
for all SKT and SGBM experiments in this report.

\subsection{2025 Ergonomics Dataset}

The current evaluation uses a single-subject ergonomics recording collected in
a controlled manufacturing-like environment.
The subject performs walking, sitting, squatting, reaching, lifting, and
interaction with furniture.
After left--right pairing by hardware frame ID, the usable synchronized camera
timeline contains 2801 frame pairs over 248.8 s, corresponding to an effective
frame rate of approximately 12.5 fps.
The left metadata contains 3015 rows and the right metadata contains 2882 rows,
so some rows are unpaired or outside the synchronized interval.

\begin{table}[H]
  \centering
  \caption{Summary of the synchronized camera timeline used for evaluation.}
  \label{tab:timeline_summary}
  \begin{tabular}{lc}
    \toprule
    Quantity & Value \\
    \midrule
    Synchronized frame pairs & 2801 \\
    Duration & 248.8 s \\
    Median frame interval & 0.0800 s \\
    Effective frame rate & 12.50 fps \\
    Left metadata rows & 3015 \\
    Right metadata rows & 2882 \\
    Rotation applied & 180\degr \\
    \bottomrule
  \end{tabular}
\end{table}

\subsection{Xsens Reference Data and Angle Representations}

The Xsens recording provides full-body motion data on an independent time axis.
Two angle representations are considered.
The first is the native Xsens joint-angle output, which follows the Xsens body
model and its internal angle definitions.
The second is the Xsens-derived geometric angle representation, where joint
angles are recomputed from Xsens 3D positions using the same three-point
formula as the vision pipelines.
The geometric representation is used as the main reference because it reduces
angle-definition mismatch between systems.

The comparison between Xsens native and Xsens-derived geometric angles is also
used as a reference-system sanity check.
For elbow flexion, the two Xsens-derived representations are very close
(about 1--1.5\degr\ MAE in the current evaluation), whereas some shoulder and
hip angles differ more strongly.
This supports the decision to report motion-space metrics in addition to
absolute angle errors.

\subsection{Cross-system Time Synchronization}

The camera timeline and Xsens timeline are aligned using an automatic offset
search.
The search first evaluates a coarse range from 0 to 30 s, then performs a fine
search around the best candidate.
Three scores are computed: a position-space score, an angle-space score, and a
motion-delta score.
The selected offset is the motion-delta optimum when sufficient samples are
available, because the motion evaluation is most sensitive to temporal
misalignment.

For the current dataset, the selected offset is 17.31 s.
The best position, angle, and motion-delta offsets are close to each other
(17.27 s, 17.28 s, and 17.31 s respectively), which indicates that a constant
offset is a reasonable alignment model for this recording.

% =============================================================
\section{Methods}
\label{sec:methods}
% =============================================================

\subsection{Overview}

This thesis compares four pose-estimation paths under a shared evaluation
protocol.
SKT is the main geometry-based stereo method.
FastSAM3D is an external monocular learning-based baseline.
Merge is an external hybrid output that combines information from FastSAM3D
and Viscando data processing.
SGBM is an alternative dense stereo path used to test whether dense disparity
can improve depth estimation at body joints.

\subsection{Sparse Keypoint Triangulation}

SKT estimates 3D body keypoints by combining 2D pose detection with calibrated
stereo geometry.
Given synchronized left and right video frames, the pipeline first applies a
2D human pose detector to each view.
The resulting COCO-format 2D keypoints are paired across cameras, filtered by
confidence and geometric consistency, and triangulated into 3D coordinates
using the calibrated projection matrices.
Joint angles are then computed from triplets of 3D keypoints.

The motivation for SKT is interpretability.
Each 3D point is linked to two observed 2D detections and a known stereo
calibration model, so the system can expose quality signals such as confidence,
epipolar error, disparity, and reprojection error.
This is valuable in industrial deployment because unreliable frames can be
detected rather than silently converted into ergonomic scores.
The main weakness is that each frame is still strongly dependent on 2D
keypoint quality; small 2D jitter can be amplified by triangulation and then
further amplified by joint-angle computation.

\subsection{FastSAM3D}

FastSAM3D is used as a learning-based monocular comparison method.
Instead of explicitly triangulating left--right keypoint pairs, FastSAM3D
produces 3D pose outputs from a single camera stream using a model-based pose
pipeline.
In this thesis, the FastSAM3D outputs are provided as TRC marker files and are
aligned to the common camera timeline before evaluation.

The motivation for including FastSAM3D is that modern learning-based pose
pipelines can implicitly encode body-shape and motion priors.
This can make the output smoother and more anatomically plausible than a
frame-independent sparse stereo reconstruction, even though the method uses
less direct geometric information.
The limitation is that monocular 3D output may not share the same absolute
scale, coordinate convention, or keypoint semantics as the stereo pipeline.
For this reason, FastSAM3D is evaluated primarily in angle and motion space,
while position-space results are interpreted cautiously.

\subsection{Merge Hybrid Output}

The Merge output is included as an external hybrid comparison.
It is evaluated with the same angle and motion metrics as SKT and FastSAM3D.
The purpose is to test whether combining Viscando processing with FastSAM3D
improves agreement with the reference signal.
In the current results, Merge does not improve the headline angle or
motion-space metrics, so it is treated as a diagnostic comparison rather than
as the main thesis method.

\subsection{Dense Stereo SGBM}

The SGBM path estimates depth from dense stereo disparity rather than sparse
keypoint triangulation.
After stereo rectification, Semi-Global Block Matching computes a disparity
map for each frame.
For each 2D keypoint location, the local disparity is converted to depth using
the stereo baseline and rectified focal length.

The motivation for SGBM is to test whether dense depth can provide a more
robust depth estimate than sparse triangulation when 2D keypoint confidence is
low.
However, dense stereo relies on local texture.
In the current dataset, dark clothing and dark furniture create low-texture
regions near important joints, causing disparity values to collapse or become
unreliable.
SGBM is therefore kept as an alternative path and diagnostic experiment, not as
the main pose-estimation method.

% =============================================================
\section{Evaluation Framework}
\label{sec:evaluation}
% =============================================================

\subsection{Angle-space Evaluation}

Angle-space evaluation measures whether the estimated posture is suitable for
ergonomic interpretation.
Eight full-body angles are evaluated: left and right shoulder, elbow, hip, and
knee flexion.
For each angle, the same three-point geometric formula is applied to the
vision methods and to the Xsens-derived reference:
\begin{equation}
  \theta = \arccos\!\left(
    \frac{(\mathbf{p}_A-\mathbf{p}_B)\cdot(\mathbf{p}_C-\mathbf{p}_B)}
         {\|\mathbf{p}_A-\mathbf{p}_B\|\;\|\mathbf{p}_C-\mathbf{p}_B\|}
  \right).
\end{equation}
The main metrics are MAE, median absolute error, bias, valid-pair ratio, and
RULA-like bin agreement for upper-body joints.

\subsection{Motion-space Evaluation}

Motion-space evaluation compares relative angle changes rather than absolute
angles.
For a given time scale $K$, the K-frame delta is defined as
\begin{equation}
  \Delta_K \theta_t = \theta_t - \theta_{t-K}.
\end{equation}
This reduces the influence of static offsets and focuses on whether two
systems describe the same motion.
The current evaluation reports $K \in \{1,6,12,25\}$, with K=6 used as the
main visualization scale because it corresponds to approximately 0.48 s at
12.5 fps.
Metrics include Pearson correlation, Spearman correlation, delta MAE, RMSE,
high-delta counts, and path ratio.

\subsection{Segment-level Evaluation}

Segment-level evaluation summarizes active movement intervals.
For each detected activity segment, range of motion (ROM), angle MAE, RULA-like
agreement, and dynamic time warping (DTW) distance are computed.
Before DTW, each segment is mean-subtracted and L2-normalized so that the
distance emphasizes motion shape rather than absolute offset or amplitude.
This makes DTW complementary to RMSE and MAE: MAE measures amplitude error,
whereas DTW measures whether the shape of the motion is similar under small
temporal deviations.

\subsection{Position-space Evaluation}

Position-space evaluation measures 3D joint distance and geometric consistency.
For stereo methods, calibrated geometry gives metric-scale coordinates, making
MPJPE and pelvis-relative comparisons meaningful after rigid alignment.
For monocular FastSAM3D output, position-space comparison is more difficult
because the output coordinate system and scale do not necessarily match the
stereo or Xsens coordinate frames.
Therefore, position-space results are used mainly to discuss the practical
advantage of stereo geometry and the coordinate/scale limitations of monocular
methods.

\subsection{Filtering Policy}

The revised evaluation uses a system-specific filtering policy.
Xsens signals are interpolated onto the shared timeline but are not smoothed
again, because they are already heavily processed internally.
Camera-based methods are evaluated after a lightweight smoothing policy
defined in the pipeline configuration.
The important methodological principle is that smoothing is applied before
delta computation; otherwise, the delta metric would measure filter artifacts
rather than the motion signal itself.

% =============================================================
\section{Results}
\label{sec:results}
% =============================================================

\subsection{Angle-space Results}

FastSAM3D achieves the lowest mean angle MAE among the evaluated vision
methods on the current dataset.
Across the eight evaluated angles, FastSAM3D obtains 15.4\degr\ mean MAE,
compared with 17.6\degr\ for SKT and 24.2\degr\ for Merge.
The result indicates that the monocular learning-based output is more stable
for angle estimation in this recording, especially for shoulder, hip, and knee
angles.
The elbow remains difficult for all vision methods.

\begin{table}[H]
  \centering
  \caption{Angle-space summary across eight evaluated joint angles. Lower MAE is better.}
  \label{tab:angle_summary}
  \small
  \begin{tabular}{lcccc}
    \toprule
    Method & Mean MAE $\downarrow$ & Median MAE $\downarrow$ &
    Valid ratio $\uparrow$ & Bin agreement $\uparrow$ \\
    \midrule
    SKT & 17.6 & 17.9 & 0.746 & 0.540 \\
    \good{FastSAM3D} & \good{15.4} & \good{14.4} & \good{0.905} & \good{0.553} \\
    Merge & 24.2 & 22.4 & 0.905 & 0.440 \\
    Xsens native check & 4.7 & 3.7 & 0.905 & 0.886 \\
    \bottomrule
  \end{tabular}
\end{table}

\subsection{Motion-space Results}

FastSAM3D also shows stronger motion agreement than SKT at K=6.
The average K=6 Pearson correlation across the eight angles is 0.64 for
FastSAM3D and 0.38 for SKT.
For the left and right elbow, which are the most important joints in the
motion-space analysis, FastSAM3D reaches Pearson correlations of 0.54 and 0.69,
whereas SKT reaches 0.39 and 0.39.
Merge has near-zero correlation in the current run, indicating that the hybrid
output does not preserve the same short-term motion pattern.

\begin{table}[H]
  \centering
  \caption{K=6 motion-delta agreement. Higher correlation is better.}
  \label{tab:motion_summary}
  \small
  \begin{tabular}{lcccc}
    \toprule
    Method & Mean $r$ $\uparrow$ & Mean $\rho$ $\uparrow$ &
    L-elbow $r$ $\uparrow$ & R-elbow $r$ $\uparrow$ \\
    \midrule
    SKT & 0.379 & 0.334 & 0.392 & 0.386 \\
    \good{FastSAM3D} & \good{0.644} & \good{0.635} & \good{0.542} & \good{0.689} \\
    Merge & 0.006 & 0.027 & 0.016 & 0.081 \\
    Xsens native check & 0.965 & 0.936 & 0.988 & 0.993 \\
    \bottomrule
  \end{tabular}
\end{table}

\begin{figure}[H]
  \centering
  \begin{subfigure}[b]{0.48\linewidth}
    \includegraphics[width=\linewidth]{scatter_skt_vs_xsensfair_leftelbow_k6.png}
    \caption{SKT vs.\ reference, left elbow.}
  \end{subfigure}
  \hfill
  \begin{subfigure}[b]{0.48\linewidth}
    \includegraphics[width=\linewidth]{scatter_fastsam3d_vs_xsensfair_leftelbow_k6.png}
    \caption{FastSAM3D vs.\ reference, left elbow.}
  \end{subfigure}
  \caption{K=6 elbow motion-delta scatter plots. FastSAM3D follows the reference trend more consistently, while SKT shows more high-delta jitter and vertical-axis accumulation.}
  \label{fig:k6_elbow_scatter}
\end{figure}

\subsection{Segment-level Results}

Segment-level ROM and DTW results reinforce the frame-level motion findings.
FastSAM3D has the lowest average ROM absolute error and the lowest DTW distance
among the vision methods, while SKT remains more variable across segments.
The Xsens native check gives the expected lower-bound reference, with small
ROM error and very low DTW distance.

\begin{table}[H]
  \centering
  \caption{Segment-level summary over detected active movement intervals. Lower values are better.}
  \label{tab:segment_summary}
  \begin{tabular}{lccc}
    \toprule
    Method & Segments & Mean ROM error (\degr) $\downarrow$ & Mean DTW distance $\downarrow$ \\
    \midrule
    SKT & 39 & 26.9 & 0.0435 \\
    \good{FastSAM3D} & 45 & \good{14.9} & \good{0.0355} \\
    Merge & 45 & 30.0 & 0.0772 \\
    Xsens native check & 45 & 4.8 & 0.0062 \\
    \bottomrule
  \end{tabular}
\end{table}

\subsection{Position-space Results}

Position-space evaluation gives a different interpretation from the angle and
motion results.
SKT achieves approximately 30.6 cm MPJPE after rigid alignment under the
current distance-evaluation protocol.
FastSAM3D has much larger direct distance to the Xsens-derived reference under
the same coordinate assumptions, because its TRC output uses a different
coordinate and scale convention.
When FastSAM3D is aligned directly to SKT over common joints, the mean
inter-method distance is approximately 60.0 cm.

\begin{table}[H]
  \centering
  \caption{Position-space diagnostic results. These values should be interpreted as coordinate/scale diagnostics rather than pure pose-quality rankings.}
  \label{tab:position_summary}
  \small
  \begin{tabularx}{\linewidth}{lccX}
    \toprule
    Comparison & Mean distance & Median distance & Notes \\
    \midrule
    SKT vs.\ reference & 30.6 cm & 20.8 cm & Metric stereo scale after rigid alignment \\
    FastSAM3D vs.\ reference & 300.3 cm & 294.8 cm & Coordinate/scale convention mismatch dominates \\
    FastSAM3D aligned to SKT & 60.0 cm & 36.3 cm & Direct inter-method skeleton comparison \\
    \bottomrule
  \end{tabularx}
\end{table}

\subsection{Cross-method Synthesis}

The three evaluation dimensions lead to a more nuanced conclusion than a
single metric would provide.
FastSAM3D is currently stronger for angle and motion agreement, suggesting that
learned body and motion priors help suppress jitter in this dataset.
SKT is weaker in short-term motion stability, but it remains valuable because
its 3D coordinates are tied to calibrated stereo geometry and because it
provides explicit quality signals.
Merge does not improve the current headline results.
SGBM shows that dense disparity is not automatically better than sparse
triangulation when the relevant body regions lack reliable texture.

% =============================================================
\section{Discussion}
\label{sec:discussion}
% =============================================================

\subsection{Why FastSAM3D Performs Better in Angle and Motion Space}

FastSAM3D likely benefits from learned priors about human body structure and
motion continuity.
Even though it uses monocular input, the model output is less dependent on
left--right keypoint matching and therefore avoids some stereo-specific jitter
failure modes.
This is particularly visible in the K=6 delta scatter plots, where SKT produces
large short-term deltas even when the reference motion is relatively stable.
Those outliers suggest that the main SKT weakness is not only random noise, but
also occasional keypoint-chain instability around the shoulder--elbow--wrist
geometry.

\subsection{Why Stereo Still Matters}

Stereo geometry remains important because it provides physical scale and
interpretable quality checks.
For real industrial deployment, knowing when the system is uncertain can be as
important as obtaining a smooth pose sequence.
SKT can expose epipolar error, reprojection error, disparity, and left--right
confidence consistency, which are directly tied to the camera geometry.
These signals can support quality gating, uncertainty reporting, or fallback
logic in a production system.

\subsection{Limitations}

The current results are based on a single subject and a single controlled
recording.
Therefore, the quantitative ranking between methods should be treated as
evidence for this dataset rather than a general conclusion about all
manufacturing environments.
The motion-space analysis is strongest for elbow motion because that joint was
investigated most deeply during the project.
The evaluation also depends on alignment, filtering, and segment-detection
parameters, although the current pipeline records these settings explicitly.
Finally, Xsens is an external comparison system rather than an absolute
physical ground truth, so absolute errors should be interpreted with this
limitation in mind.

\subsection{Implications for Real-time Deployment}

The real-time deployment path should first focus on making the existing SKT
pipeline efficient on GPU hardware.
The immediate engineering steps are to export the 2D pose model to TensorRT,
use GPU-accelerated video decoding and preprocessing, and keep triangulation
and angle computation lightweight enough for streaming.
If this baseline runs reliably, transformer-based temporal priors or
FastSAM3D-like model outputs can be explored as optional stabilizers.
FastSAM3D's current runtime on GPU hardware can serve as a practical benchmark
for the expected performance level of learning-based pose methods.

% =============================================================
\section{Conclusion and Future Work}
\label{sec:conclusion}
% =============================================================

\subsection{Conclusion}

This thesis draft presents a multi-dimensional evaluation of 3D pose estimation
for manufacturing-oriented ergonomic assessment.
The main finding is that method quality depends strongly on the evaluation
dimension.
FastSAM3D currently provides better angle-space and motion-space agreement on
the evaluated recording, while SKT provides more interpretable metric 3D
geometry and quality signals.
This suggests that future ergonomic monitoring systems may benefit from a
hybrid view: learning-based models for stable body-pose priors, and stereo
geometry for scale, quality control, and deployment transparency.

\subsection{Future Work}

The next step is validation on additional recordings and subjects.
This is necessary to test whether the current FastSAM3D advantage generalizes
and whether SKT failure modes are consistent across operators, clothing,
lighting, and camera distances.
The second step is real-time GPU deployment of the SKT pipeline, followed by
optional temporal-model integration if the computational budget allows.
The third step is to extend the motion-space evaluation beyond the elbow joint
and connect pose-estimation outputs more directly to full ergonomic risk
assessment.

% =============================================================
\appendix
\clearpage
% =============================================================

\section{Implementation and Reproducibility Notes}

The standalone evaluation workflow is implemented in \texttt{00\_pose\_pipeline/}.
For the current dataset, the main reproduction command is:
\begin{footnotesize}
\begin{verbatim}
/opt/anaconda3/envs/pose/bin/python \
  00_pose_pipeline/src/run_pipeline.py \
  --config \
  00_pose_pipeline/configs/current_2025_ergonomics.yaml \
  --stages \
  validate,offset,angle,motion,segment,scatter
\end{verbatim}
\end{footnotesize}

The configuration records dataset-specific choices such as 180\degr\ video
rotation, hardware-frame-ID synchronization, timestamp parsing, camera
calibration path, smoothing settings, and the automatic offset-search range.
Generated outputs are stored under \texttt{00\_pose\_pipeline/runs/}.

\section{Claim-evidence Self-review}

\begin{table}[H]
  \centering
  \caption{Claim-evidence map for the current thesis draft.}
  \label{tab:claim_evidence}
  \small
  \begin{tabularx}{\linewidth}{p{4.0cm}Xp{1.9cm}}
    \toprule
    Claim & Evidence used in this draft & Status \\
    \midrule
    FastSAM3D is stronger than SKT in angle-space on the current dataset.
    & Table~\ref{tab:angle_summary}: 15.4\degr\ mean MAE vs.\ 17.6\degr.
    & Supported \\
    FastSAM3D is stronger than SKT in K=6 motion-space agreement.
    & Table~\ref{tab:motion_summary}: mean Pearson 0.644 vs.\ 0.379; elbow-specific correlations also higher.
    & Supported \\
    SKT retains a position-space and interpretability advantage.
    & Table~\ref{tab:position_summary}; SKT uses calibrated stereo scale and exposes geometric quality signals.
    & Supported \\
    Results generalize to manufacturing environments broadly.
    & Only one controlled single-subject dataset is evaluated.
    & Needs evidence \\
    Full RULA automation is solved.
    & Only joint-angle and RULA-like bin agreement are evaluated.
    & Not claimed \\
    \bottomrule
  \end{tabularx}
\end{table}

\begin{thebibliography}{9}

\bibitem{kabsch1976solution}
W.~Kabsch,
\textit{A solution for the best rotation to relate two sets of vectors},
Acta Crystallographica, Section A, vol.~32, pp.~922--923, 1976.

\bibitem{hirschmuller2007sgbm}
H.~Hirschm{\"u}ller,
\textit{Stereo processing by semiglobal matching and mutual information},
IEEE Transactions on Pattern Analysis and Machine Intelligence,
vol.~30, no.~2, pp.~328--341, 2007.

\end{thebibliography}

\end{document}
